% Options for packages loaded elsewhere
\PassOptionsToPackage{unicode}{hyperref}
\PassOptionsToPackage{hyphens}{url}
\documentclass[
]{article}
\usepackage{xcolor}
\usepackage{amsmath,amssymb}
\setcounter{secnumdepth}{-\maxdimen} % remove section numbering
\usepackage{iftex}
\ifPDFTeX
  \usepackage[T1]{fontenc}
  \usepackage[utf8]{inputenc}
  \usepackage{textcomp} % provide euro and other symbols
\else % if luatex or xetex
  \usepackage{unicode-math} % this also loads fontspec
  \defaultfontfeatures{Scale=MatchLowercase}
  \defaultfontfeatures[\rmfamily]{Ligatures=TeX,Scale=1}
\fi
\usepackage{lmodern}
\ifPDFTeX\else
  % xetex/luatex font selection
\fi
% Use upquote if available, for straight quotes in verbatim environments
\IfFileExists{upquote.sty}{\usepackage{upquote}}{}
\IfFileExists{microtype.sty}{% use microtype if available
  \usepackage[]{microtype}
  \UseMicrotypeSet[protrusion]{basicmath} % disable protrusion for tt fonts
}{}
\makeatletter
\@ifundefined{KOMAClassName}{% if non-KOMA class
  \IfFileExists{parskip.sty}{%
    \usepackage{parskip}
  }{% else
    \setlength{\parindent}{0pt}
    \setlength{\parskip}{6pt plus 2pt minus 1pt}}
}{% if KOMA class
  \KOMAoptions{parskip=half}}
\makeatother
\usepackage{longtable,booktabs,array}
\newcounter{none} % for unnumbered tables
\usepackage{calc} % for calculating minipage widths
% Correct order of tables after \paragraph or \subparagraph
\usepackage{etoolbox}
\makeatletter
\patchcmd\longtable{\par}{\if@noskipsec\mbox{}\fi\par}{}{}
\makeatother
% Allow footnotes in longtable head/foot
\IfFileExists{footnotehyper.sty}{\usepackage{footnotehyper}}{\usepackage{footnote}}
\makesavenoteenv{longtable}
\setlength{\emergencystretch}{3em} % prevent overfull lines
\providecommand{\tightlist}{%
  \setlength{\itemsep}{0pt}\setlength{\parskip}{0pt}}
\usepackage{bookmark}
\IfFileExists{xurl.sty}{\usepackage{xurl}}{} % add URL line breaks if available
\urlstyle{same}
\hypersetup{
  hidelinks,
  pdfcreator={LaTeX via pandoc}}

\author{}
\date{}

\begin{document}

\section{The Architectural Foundations of Navier--Stokes in Event
Density}\label{the-architectural-foundations-of-navierstokes-in-event-density}

\subsection{Form Derivation, Clay-Relevance Decomposition, and the
Structural Status of
Turbulence}\label{form-derivation-clay-relevance-decomposition-and-the-structural-status-of-turbulence}

\textbf{Allen Proxmire} \textbf{Collaborator:} Claude (AI collaborator)
\textbf{April 2026} \textbf{Series:} Event-Density Foundational Theorems
--- Navier--Stokes Synthesis

\begin{center}\rule{0.5\linewidth}{0.5pt}\end{center}

\subsection{Abstract}\label{abstract}

Classical fluid mechanics describes the Navier--Stokes equations as a
phenomenological set of conservation laws without an architectural
explanation of why the equations take their specific form, why 2D NS is
globally smooth while 3D NS smoothness remains the open Clay problem, or
why turbulence lacks a canonical structural template across systems. We
present a unified architectural account of these questions within the
Event Density (ED) program, integrating five closed structural arcs.
\textbf{NS-1} establishes that ED's canonical PDE forces D = 3+1 via a
Path B-strong argument combining architectural d ≥ 3 (T20 mechanism
degeneracy at d ≤ 2), architectural-suggestive d ≤ 3 (Polya recurrence +
concentration of measure + decoupling-surface degeneration), and
empirical-consistency closure (T19/T20 outputs match observed gravity at
d = 3). \textbf{NS-2} derives the standard Newtonian-fluid NS form from
substrate primitives via two concordant routes --- chain-substrate
coarse-graining and partial vector-extension of the canonical PDE ---
with advection, pressure, and incompressibility catalogued as
fluid-mechanical additions not native to ED canonical channels.
\textbf{NS-3} reaches an Intermediate Path C verdict: ED contains a real
Clay-NS-relevant regularizing mechanism (R1: form-FORCED
\(-\kappa\mu_{V1}\ell_P^2 \nabla^4 v\) stabilization arising from V1's
finite-width vacuum kernel), with quantitative competition against
destabilizing super-Burnett terms INHERITED on both sides.
\textbf{NS-Smoothness} formalizes the Clay-relevance decomposition: R1
produces strictly monotone gradient-norm Lyapunov decay in ED-only NS
(without advection), while the advective vortex-stretching term
\(\int\omega\cdot S\omega\,dV\) is the unique indefinite-sign
contribution to \(dL/dt\) in full 3D NS. \textbf{NS-Turb} closes the P7
↔ turbulence-cascade mapping at H1 trivial / H2 partial / H3 fail --- no
architectural template for developed-cascade turbulence. Three
independent program-level analyses converge on advection as structurally
non-ED at architectural (NS-2.08), dynamical (NS-Smooth-3), and spectral
(NS-Turb-4) levels. The synthesis yields a unified structural picture:
ED supplies real architectural regularization and canonical NS
form-derivation, but advection --- the unique obstruction to ED-style
smoothness --- is structurally non-ED at three independent levels. Event
Density explains why the NS form appears, why 2D NS is globally smooth
(vortex-stretching vanishes identically), and why 3D NS smoothness is
structurally hard (the obstruction lies outside the canonical
regularizing architecture). It does not resolve whether 3D NS blows up
at finite time, predict critical Reynolds numbers, or supply a
turbulence-cascade architectural template. This paper provides the first
unified architectural account of Navier--Stokes within the Event Density
program, complementing the empirical P4-NN rheology paper that covers
ED's reach into non-Newtonian fluid mechanics.

\begin{center}\rule{0.5\linewidth}{0.5pt}\end{center}

\subsection{1. Introduction}\label{introduction}

\subsubsection{1.1 Motivation}\label{motivation}

The Navier--Stokes equations are the empirically successful description
of viscous incompressible fluid flow, validated across a vast range of
laboratory and engineering regimes. Their phenomenological success
contrasts sharply with the structural questions they leave open. Why
does the velocity field obey precisely the equation it obeys, and not
some other equation with similar dimensional content? Why does the same
equation in two spatial dimensions admit global smooth solutions for
arbitrary finite-energy data, while in three spatial dimensions the
analogous question --- the Clay Millennium smoothness problem --- has
remained open for nine decades? Why does turbulence, the developed
nonlinear regime of 3D NS, lack a single canonical structural account
that crosses systems and Reynolds numbers? Standard fluid mechanics
treats these questions as either resolved by observation (the equation
works), as deferred to mathematics (Clay-NS), or as phenomenological
complexity (turbulence). What is missing is an architectural account of
the equation itself.

This paper presents such an account within the Event Density (ED)
program, an architectural framework whose canonical scalar partial
differential equation derives from substrate-level commitments rather
than fluid-mechanical postulates. The relationship between ED and
Navier--Stokes turns out to be neither one of identity nor of
irrelevance. ED reproduces the \emph{viscous} content of NS at
architectural level, derives the \emph{dimensional structure} (D = 3+1)
from substrate-level arguments, and supplies a \emph{regularizing
mechanism} that would close the Clay smoothness question were it not for
a structural feature ED does not natively absorb --- the advective
convective derivative.

\subsubsection{1.2 Event Density as an architectural
framework}\label{event-density-as-an-architectural-framework}

Event Density is a relational substrate ontology. Its primitives include
discrete micro-events, channels along which micro-events propagate,
chains as worldline-class persistent structures, decoupling surfaces as
participation thresholds, and primitive-level commitment
irreversibility. From these substrate primitives plus seven structural
canon principles (P1--P7), ED derives a canonical scalar nonlinear PDE
that has been shown to generate, in three nominally separate domains,
the foundational structures of physics: the four postulates of
non-relativistic quantum mechanics, the substrate-derived form of
Newton's gravitational law together with the MOND transition
acceleration and the slope-4 baryonic Tully--Fisher relation, and a
universal degenerate-mobility law that fits soft-matter diffusion data
across ten chemically unrelated systems with \(R^2 > 0.986\) at common
exponent \(\beta = 1.72 \pm 0.37\).

The ED program is summarized in {[}\emph{Event Density: One Substrate,
Three Domains}{]} (Proxmire 2026); the present paper extends the
program's structural reach into mainstream classical fluid mechanics,
focused specifically on the Navier--Stokes equations and the Clay
smoothness problem.

\subsubsection{1.3 The five closed Navier--Stokes
arcs}\label{the-five-closed-navierstokes-arcs}

The structural content of this paper integrates five closed program
arcs:

\begin{itemize}
\tightlist
\item
  \textbf{NS-1} (Path B-strong dimensional forcing): ED's canonical PDE
  forces D = 3+1 via three concordant lines.
\item
  \textbf{NS-2} (substrate derivation of NS form): standard
  Newtonian-fluid NS form derivable via two distinct substrate routes.
\item
  \textbf{NS-3} (Intermediate Path C): ED contains a real
  Clay-NS-relevant regularizing mechanism with quantitative competition
  INHERITED.
\item
  \textbf{NS-Smoothness} (Clay-relevance decomposition): formal R1 +
  advection-obstruction decomposition.
\item
  \textbf{NS-Turb} (turbulence cascade): no architectural template for
  developed-cascade turbulence.
\end{itemize}

Each arc is structurally complete; this paper consolidates and
integrates them.

\subsubsection{1.4 Scope}\label{scope}

The paper presents the architectural-decompositional content of ED's
relationship to Navier--Stokes. It does not solve the Clay-NS smoothness
problem, predict critical Reynolds numbers for any specific transition,
supply turbulence closure models (LES, RANS, sub-grid), or propose
architectural extensions to absorb advection canonically. It catalogues
what is delivered by the existing canon and what lies beyond.

The empirical content of ED's reach into non-Newtonian fluid mechanics
(jamming, discontinuous shear-thickening, Cross/Carreau shear-thinning,
Maxwell-class viscoelasticity) is the subject of a companion paper,
\emph{ED Mobility Saturation Predicts Non-Newtonian Rheology} (the P4-NN
paper). The two papers together cover the program's published-grade
content on classical fluids.

\subsubsection{1.5 Contributions}\label{contributions}

\begin{itemize}
\tightlist
\item
  A unified architectural account of why the Navier--Stokes form
  appears, derived from substrate primitives via two concordant routes.
\item
  A formal Clay-relevance decomposition: ED's R1 mechanism (positive
  side) + advective vortex-stretching (negative side, structurally
  non-ED).
\item
  A three-angle convergence theorem: advection is identified as
  structurally non-ED at architectural, dynamical, and spectral levels
  independently.
\item
  A structural account of the 2D-versus-3D smoothness asymmetry:
  vortex-stretching, the obstruction to ED-style monotone gradient-norm
  decay, vanishes identically in 2D and is sign-indefinite in 3D.
\item
  An honest catalogue of what ED explains (form derivation, dimensional
  forcing, structural decomposition of Clay difficulty) and what it does
  not (Clay-NS resolution, Reynolds-number predictions,
  turbulence-cascade template).
\end{itemize}

\subsubsection{1.6 Paper organization}\label{paper-organization}

Section 2 lays out the ED architectural background. Section 3 derives
the Navier--Stokes form from ED via the two concordant substrate routes.
Section 4 introduces the R1 regularizing mechanism. Section 5 analyzes
the advection obstruction. Section 6 establishes the three-angle
convergence on advection-as-non-ED. Section 7 presents the Intermediate
Path C Clay-relevance verdict. Section 8 closes the turbulence question.
Section 9 synthesizes the unified picture. Section 10 concludes.
Appendices A--C provide key equations, closed-arc summaries, and
architectural canon references.

\begin{center}\rule{0.5\linewidth}{0.5pt}\end{center}

\subsection{2. Architectural Background}\label{architectural-background}

\subsubsection{2.1 The ED canonical PDE}\label{the-ed-canonical-pde}

The canonical Event Density partial differential equation, established
in the program's foundational papers, takes the form

\[\partial_t \rho \;=\; D \cdot F[\rho] \;+\; H \cdot v, \qquad D + H = 1, \qquad D, H \in [0, 1],\]

\[\dot v \;=\; \tau^{-1}\bigl(\bar F(\rho) - \zeta\, v\bigr),\]

with the operator

\[F[\rho] \;=\; M(\rho)\,\nabla^2 \rho \;+\; M'(\rho)\,|\nabla\rho|^2 \;-\; P(\rho).\]

Three constitutive channels --- mobility (\(M\)), penalty (\(P\)), and
global participation (\(v\)) --- are jointly necessary and sufficient to
satisfy seven structural principles P1--P7 that define the ED
equivalence class. The canonical scalar PDE acts on a bounded density
field \(\rho(\mathbf{x}, t)\). A vector-extension framework, articulated
in Architectural\_Canon\_Vector\_Extension, lifts the architectural
principles to vector and tensor field structures while preserving the
architectural canon's content.

\subsubsection{2.2 P1--P7: the Architectural
Canon}\label{p1p7-the-architectural-canon}

The seven canon principles are stated in compressed form below. Detailed
treatments are in the foundational ED literature.

\begin{itemize}
\tightlist
\item
  \textbf{P1 --- Operator structure.}
  \(F[\rho] = M(\rho)\nabla^2\rho + M'(\rho)|\nabla\rho|^2 - P(\rho)\)
  is the unique three-term form satisfying locality, isotropy, and the
  dissipative-structure constraints.
\item
  \textbf{P2 --- Channel complementarity.}
  \(\partial_t\rho = D \cdot F[\rho] + H \cdot v\) with \(D + H = 1\),
  splitting dynamics into direct (substrate operator) and mediated
  (participation feedback) channels.
\item
  \textbf{P3 --- Penalty equilibrium.} The penalty function \(P(\rho)\)
  has a unique zero at \(\rho = \rho^*\), with \(P'(\rho^*) > 0\),
  ensuring monostable equilibrium structure.
\item
  \textbf{P4 --- Mobility capacity bound.} The mobility function
  \(M(\rho)\) satisfies \(M(\rho_\mathrm{max}) = 0\) with
  \(M(\rho) > 0\) for \(\rho < \rho_\mathrm{max}\), encoding finite
  packing capacity.
\item
  \textbf{P5 --- Participation feedback.} The global participation
  variable \(v(t)\) obeys
  \(\dot v = \tau^{-1}(\bar F(\rho) - \zeta v)\), supplying memory and
  oscillatory structure.
\item
  \textbf{P6 --- Damping discriminant.} A canonical damping discriminant
  determines underdamped versus overdamped regimes at
  \(D_\mathrm{crit} \approx 0.896\) (canonical \(\zeta = 1/4\)).
\item
  \textbf{P7 --- Nonlinear triad coupling.} The nonlinearity
  \(M'(\rho)|\nabla\rho|^2\) generates harmonic content at the
  few-percent level (canonical 3--6\% at \(k = 3\) from \(k = 1\)
  driving) under multiplicative perturbations.
\end{itemize}

\subsubsection{2.3 Vector extension}\label{vector-extension}

The canonical PDE acts on a scalar density. The Architectural Canon
Vector Extension establishes that P1--P7 are field-type-agnostic at the
architectural level: vector and tensor field PDEs satisfying P1--P7
component-wise are architecturally ED, even when they violate the
concrete-PDE-level constraint C5 (single scalar field) of the canonical
exemplar. A three-tier classification distinguishes: canonical (scalar,
satisfying both C-level and P-level constraints), fully ED-architectural
(vector/tensor, satisfying P-level), and partially ED-architectural
(containing P-level content plus structural additions not native to the
canon).

This framework will be load-bearing in Section 3, where we treat the
Navier--Stokes velocity field as a candidate for
partial-ED-architectural status.

\subsubsection{2.4 Form-FORCED versus value-INHERITED
methodology}\label{form-forced-versus-value-inherited-methodology}

A persistent distinction throughout the ED program separates:

\begin{itemize}
\tightlist
\item
  \textbf{Form-FORCED} content --- functional forms derivable from
  architectural principles and admitting no alternative within the
  canon.
\item
  \textbf{Value-INHERITED} content --- specific numerical parameters
  that the canon does not determine and that are inherited from
  material-specific physics or experimental input.
\end{itemize}

This distinction is methodological: the ED program does not claim to
predict every numerical value in physics from primitives alone. It
claims that the \emph{forms} of physical laws --- the functional
structures of equations and the qualitative existence of phenomena ---
follow as theorems from substrate commitments. Specific values fall
under empirical or material-specific physics. The Navier--Stokes program
preserves this discipline throughout.

\begin{center}\rule{0.5\linewidth}{0.5pt}\end{center}

\subsection{3. Deriving Navier--Stokes from Event
Density}\label{deriving-navierstokes-from-event-density}

\subsubsection{3.1 Path B-strong: D = 3+1
forced}\label{path-b-strong-d-31-forced}

The first structural question for any architectural derivation of
Navier--Stokes is dimensional: why three spatial dimensions plus one
time dimension? Standard fluid mechanics takes D = 3+1 as observational
input. NS-1 closes this question via a Path B-strong argument that
combines three concordant lines:

\textbf{Architectural lower bound (\(d \ge 3\)).} The substrate-gravity
arc's transition-acceleration mechanism (T20) requires the cosmic
decoupling surface to project anisotropically onto an accelerating chain
via the dipole spherical harmonic \((\ell = 1, m = 0)\). This projection
geometry degenerates at \(d \le 2\): the cosmic horizon's \(S^{d-1}\)
has no separate polar-azimuthal decomposition for \(d - 1 \le 1\), and
the canonical 2π azimuthal periodicity that yields
\(a_0 = c \cdot H_0/(2\pi)\) has no analog. ED cannot host its own
substrate-gravity arc at \(d \le 2\) --- an internal-consistency forcing
of \(d \ge 3\).

\textbf{Architectural-suggestive upper bound (\(d \le 3\)).} Three
independent primitive-anchored mechanisms converge: (i) concentration of
measure on participation neighborhoods (in high \(d\), near/far
discrimination collapses); (ii) Polya recurrence/transience boundary
(random walks on \(\mathbb{Z}^d\) are recurrent for \(d \le 2\),
transient for \(d \ge 3\), with the \(d = 3\) marginal regime supporting
cross-chain participation overlap that V1 sourcing requires); (iii)
ED-06 decoupling-surface boundary/bulk degeneration in high-dimensional
balls. Polya is the load-bearing mechanism; concentration and
decoupling-degeneration are concordant supporting arguments. Each
carries an identification gap (chain-dynamics-as-diffusion at
coarse-grained scale; participation-bandwidth as probability measure;
ED-06 ontology as Euclidean-ball geometry). The combined argument does
not rise to a strict-architectural theorem but is structurally
suggestive.

\textbf{Empirical-consistency closure (\(d \le 3\)).} T19's
holographic-bound derivation and T20's dipole-mode mechanism produce, in
arbitrary \(d\)-spatial, gravity laws of form \(a \propto M/r^{d-1}\)
and prefactor structures dependent on \(S^{d-1}\) harmonic
decomposition. Empirical Newton (1/r²) and empirical \(a_0\)
(\(\sim 1.2 \times 10^{-10}\,\mathrm{m/s^2}\)) match the
substrate-gravity outputs only at \(d = 3\). The substrate primitives
are d-agnostic at primitive-statement level; the d = 3 specification
enters via empirical match to observed gravity.

The aggregate verdict (Path B-strong) closes B2 at multi-layer
architectural plus empirical-consistency level. D = 3+1 is forced by
ED's structural framework when taken in conjunction with observed
gravity. The forcing is multi-source, layer-attributed, and structurally
honest: the seven PDE-uniqueness axioms alone do not force D = 3+1, but
ED's full primitive-plus-substrate-plus-empirical framework does.

\subsubsection{3.2 Substrate-level derivation of the NS
form}\label{substrate-level-derivation-of-the-ns-form}

NS-2 derives the standard Newtonian-fluid Navier--Stokes form from
substrate primitives via two methodologically distinct but structurally
concordant routes.

\textbf{Chain-substrate route (NS-2.01--NS-2.07).} A coarse-graining
cell of radius \(R_\mathrm{cg}\) is defined, satisfying three
constraints simultaneously: chain-discreteness suppression
(\(R_\mathrm{cg}\) much larger than substrate UV cutoff \(\ell_P\)),
hydrodynamic regime (\(R_\mathrm{cg}\) much larger than mean free path
\(\lambda_\mathrm{mfp}\)), and field-structure preservation
(\(R_\mathrm{cg}\) much smaller than flow scale \(L_\mathrm{flow}\)).
Within this scaling window, four substrate conserved quantities are
identified --- chain count, chain mass (from Arc M), chain momentum, and
chain bandwidth content (energy-class) --- and their fluxes through the
coarse-graining cell boundary are computed via 2-sphere boundary
integration parallel to T19's holographic-bound mechanism. The resulting
flux forms

\[\mathbf{J}_\rho = \rho \mathbf{v}, \qquad \Pi_{ij} = \rho v_i v_j + \tau_{ij}, \qquad \mathbf{J}_e = e \mathbf{v} + \mathbf{v} \cdot \tau + \mathbf{Q}\]

feed directly into the continuity, momentum-balance, and energy
equations. The stress tensor \(\tau_{ij}\) decomposes into pressure plus
Newtonian viscous deviatoric plus an ED-specific residual that is
structurally identified at substrate level but numerically negligible at
NS scales. The result is the standard Newtonian-fluid Navier--Stokes
form.

\textbf{ED-PDE-direct route (NS-2.08).} Independently, the canonical
scalar PDE's three channels are mapped onto velocity-field components
via the partial vector-extension framework. The mobility/gradient
channel applied component-wise to velocity components yields the viscous
term \(\mu \nabla^2 v_i\), with the canonical \(D \cdot M\) identified
as the kinematic viscosity \(\nu\). The compositional rule's
gradient-penalty content corresponds to the same viscous structure
viewed at the configuration-space layer.

Both routes produce the same Newtonian-fluid form. The chain-substrate
route explains coefficient origin (viscosity emerges from kinetic +
cross-chain V5 + V1-mediated contributions, with values inherited at
substrate level); the ED-PDE-direct route explains architectural form
(viscous content is canonical ED architecture; pressure / advection /
incompressibility are flagged as fluid-mechanical-additions).

\subsubsection{3.3 Partial vector extension and the NS
form}\label{partial-vector-extension-and-the-ns-form}

The vector-extension framework permits NS to be classified within the
three-tier scheme of Section 2.3. NS satisfies the architectural canon's
content for its viscous component: the mobility channel applied
component-wise to velocity gives Newtonian viscous diffusion as a clean
P1-class content, and the principles P2--P7 generalize component-wise
without violation. NS does not satisfy the canon's content for its
pressure, advective, and incompressibility content: these are structural
features of fluid kinematics that lie outside the canonical P1--P7
framework.

NS is therefore \textbf{partially ED-architectural} in the three-tier
classification: its viscous content is canonical ED vector-extended; its
pressure, advection, and incompressibility content is
fluid-mechanical-specific structural addition that the canon does not
natively absorb.

\subsubsection{3.4 Identification of advection, pressure, and
incompressibility as fluid-mechanical
additions}\label{identification-of-advection-pressure-and-incompressibility-as-fluid-mechanical-additions}

The catalogue of fluid-mechanical additions, established in NS-2.08,
comprises three structural features:

\begin{itemize}
\item
  \textbf{Pressure as Lagrange multiplier.} In incompressible NS,
  pressure is a Lagrange multiplier enforcing
  \(\nabla \cdot \mathbf{v} = 0\). It is not derived from any ED
  canonical channel; it is a fluid-mechanical structural device required
  for coherence with the holonomic constraint of incompressibility.
\item
  \textbf{Advective convective derivative.} The term
  \((\mathbf{v} \cdot \nabla)\mathbf{v}\) is a kinematic coupling
  between velocity components arising from the fact that velocity both
  \emph{is} the advecting flow and \emph{is being} advected. It has no
  analog in any P-level canonical channel: it is symmetric in neither
  index pattern nor scaling structure with any P1-class operator
  content.
\item
  \textbf{Incompressibility constraint.} \(\nabla \cdot \mathbf{v} = 0\)
  is a holonomic constraint not derivable from any P-level canonical
  principle. It enters as a fluid-mechanical commitment characterizing
  the regime (low-Mach-number limit of compressible NS) rather than as a
  substrate-derived structural fact.
\end{itemize}

These three additions are structurally significant for the rest of the
paper: pressure and incompressibility, despite being
fluid-mechanical-specific, contribute zero or vanish appropriately in
the gradient-norm Lyapunov analyses of Sections 4 and 5; advection alone
produces the Clay-NS-relevant obstruction.

\subsubsection{3.5 Concordance of the two derivation
routes}\label{concordance-of-the-two-derivation-routes}

The two derivation routes --- chain-substrate kinetic-theory-class
coarse-graining and ED-PDE-direct vector-extension --- produce the same
Newtonian-fluid Navier--Stokes form. The concordance is non-trivial: the
routes operate at different methodological layers (substrate-level
statistical coarse-graining versus architectural-canon component-wise
application) and use disjoint analytical machinery, yet converge on
identical equations. This concordance constitutes structural evidence
that ED's content for the viscous part of NS is not an artifact of any
single derivation framework. The remaining structural additions
(pressure, advection, incompressibility) appear identically in both
routes as fluid-mechanical-specific content not derivable from substrate
primitives or from canon principles.

\begin{center}\rule{0.5\linewidth}{0.5pt}\end{center}

\subsection{4. The R1 Mechanism}\label{the-r1-mechanism}

\subsubsection{4.1 Origin of R1 from V1's finite-width vacuum
kernel}\label{origin-of-r1-from-v1s-finite-width-vacuum-kernel}

Theorem N1 of the Event Density program establishes the existence of a
vacuum response kernel \(K_\mathrm{vac}(x, x')\) that is Lorentz-scalar,
finite-width on the substrate scale \(\ell_\mathrm{ED}\), and
sub-power-law-2 decaying. Combined with the Newton-recovery argument in
T19 (which forces \(\ell_\mathrm{ED} = \ell_P\), the Planck length),
this kernel becomes the structural object on which the R1 mechanism is
built.

When the substrate-level vacuum response kernel is coarse-grained to
fluid-mechanical scales via multi-scale expansion in the small parameter
\(\varepsilon = \ell_P / L_\mathrm{flow}\), the leading-order
substrate-cutoff correction to the continuum Navier--Stokes momentum
equation appears as a higher-derivative regularization term:

\[\rho \frac{D v_i}{Dt} = -\partial_i p + \mu \nabla^2 v_i \;-\; \kappa \mu_{V1} \ell_P^2 \nabla^4 v_i + \rho f_i^\mathrm{ext},\]

with coefficient \(\kappa\) form-FORCED positive by V1 being a positive
smoothing kernel (its Fourier transform
\(\hat K(k) = K_0[1 - \alpha (k\ell_P)^2 + \cdots]\) has \(\alpha > 0\)
for all positive monotonically-decreasing-in-\(|k|\) kernels). The
magnitude \(\mu_{V1}\) is inherited via V1's specific G-function
profile.

\subsubsection{\texorpdfstring{4.2 Form-FORCED \(\nabla^4 v\)
stabilization}{4.2 Form-FORCED \textbackslash nabla\^{}4 v stabilization}}\label{form-forced-nabla4-v-stabilization}

The R1 term \(-\kappa \mu_{V1} \ell_P^2 \nabla^4 v_i\) is the structural
shadow of ED's substrate UV cutoff in the continuum momentum equation.
Three structural facts characterize it:

\begin{enumerate}
\def\labelenumi{\arabic{enumi}.}
\item
  \textbf{Existence is form-FORCED.} The term arises from V1's
  finite-width substrate-level kernel under multi-scale expansion. It is
  not optional: any continuum NS theory respecting ED's canonical V1
  kernel structure includes the term.
\item
  \textbf{Sign is form-FORCED positive.} The \(\kappa > 0\) sign is
  forced by V1's positive-smoothing-kernel structure. The term is
  dissipative (energy-removing in the momentum equation;
  gradient-removing in the gradient-norm Lyapunov).
\item
  \textbf{Magnitude is value-INHERITED.} The coefficient
  \(\kappa \mu_{V1}\) depends on V1's specific G-function profile, which
  is inherited from material-specific substrate physics rather than
  fixed by the canon at primitive level.
\end{enumerate}

The term is suppressed by the ratio \(\ell_P^2 / L^2 \le 10^{-60}\) at
laboratory NS scales but activates at substrate-approaching gradient
scales. This is the regime where the R1 mechanism becomes Clay-relevant:
in any putative finite-time blow-up trajectory, gradients reach scales
where R1 dominates the dissipative content of the equation.

\subsubsection{4.3 Gradient-norm Lyapunov
structure}\label{gradient-norm-lyapunov-structure}

Define the gradient-norm Lyapunov functional

\[L(t) = \frac{1}{2} \|\nabla \mathbf{v}(t)\|_2^2 = \frac{1}{2} \int_{\mathbb{R}^3} \partial_j v_i \, \partial_j v_i \, dV.\]

This is enstrophy-class --- directly related to enstrophy
\(\frac{1}{2}\int|\boldsymbol{\omega}|^2 dV\) via integration by parts
and incompressibility. By the Beale--Kato--Majda regularity criterion
combined with standard energy methods, monotonic control of this
quantity in 3D would imply global smooth solutions for finite-energy
initial data.

For the \emph{ED-only NS} equation --- the counterfactual obtained from
full NS by removing the advective convective derivative ---

\[\rho\,\partial_t v_i = \mu\nabla^2 v_i - \kappa\mu_{V1}\ell_P^2\nabla^4 v_i - \partial_i p, \qquad \nabla\cdot\mathbf{v} = 0,\]

the gradient-norm Lyapunov derivative computes as

\[\frac{dL}{dt}\bigg|_\text{ED-only NS} = -\nu \int |\nabla^2 v|^2 \, dV \;-\; \kappa \mu_{V1} \ell_P^2 \int |\nabla^3 v|^2 \, dV \;\le\; 0,\]

with viscous diffusion contributing the standard
\(-\nu \|\nabla^2 v\|_2^2\) term and R1 contributing the additional
\(-\kappa\mu_{V1}\ell_P^2 \|\nabla^3 v\|_2^2\) term via four
integrations by parts on the \(\nabla^4\) operator. Pressure contributes
zero by integration by parts combined with incompressibility.

Both contributions are manifestly non-positive. There are no positive
terms in the ED-only NS gradient-norm Lyapunov derivative; the gradient
norm is strictly monotonically decreasing along trajectories.

\subsubsection{4.4 ED-only NS global
smoothness}\label{ed-only-ns-global-smoothness}

The ED-only NS equation with \(\kappa > 0\) is a
higher-derivative-regularized parabolic equation. Standard
parabolic-regularity theory for this class (Lions 1969 and downstream
literature) gives global smooth solutions for smooth, finite-energy
initial data on \(\mathbb{R}^3\). Specifically:

\begin{enumerate}
\def\labelenumi{\arabic{enumi}.}
\tightlist
\item
  Local well-posedness in \(H^s\) for \(s\) sufficiently large.
\item
  Global \(H^s\) bound from the gradient-norm Lyapunov decay combined
  with the standard \(L^2\) energy bound.
\item
  Bootstrap to higher regularity via repeated parabolic-regularity
  application; the \(\nabla^4\) regularization makes each step strictly
  stronger than standard NS.
\end{enumerate}

This argument is canonical for higher-derivative-regularized parabolic
equations and parallels global-smoothness results for the
Navier--Stokes--Burgers equation. In our case, with the advective term
entirely absent from ED-only NS, the smoothness result is \emph{a
fortiori} --- the parabolic regularity machinery applies without the
obstructive nonlinearity.

This establishes the \textbf{positive side of the Clay-relevance
decomposition}: ED's architectural content alone, without the
fluid-mechanical-addition advection, is sufficient for global smoothness
on \(\mathbb{R}^3\). The structural significance is that ED supplies a
real, canon-level smoothing mechanism that would close the smoothness
question in a fluid lacking advection.

\begin{center}\rule{0.5\linewidth}{0.5pt}\end{center}

\subsection{5. The Advection
Obstruction}\label{the-advection-obstruction}

\subsubsection{5.1 Full NS + R1 gradient-norm
derivative}\label{full-ns-r1-gradient-norm-derivative}

When the advective convective derivative is restored to ED-only NS, full
incompressible Navier--Stokes augmented with the R1 stabilization is
recovered:

\[\rho\,\partial_t v_i + \rho\,(v \cdot \nabla)v_i \;=\; \mu\nabla^2 v_i - \kappa\mu_{V1}\ell_P^2\nabla^4 v_i - \partial_i p, \qquad \nabla \cdot \mathbf{v} = 0.\]

Computing \(dL/dt\) on this equation, three of four contributions are
unchanged from the ED-only case: viscous diffusion contributes
\(-\nu \|\nabla^2 v\|_2^2 \le 0\); R1 contributes
\(-\kappa\mu_{V1}\ell_P^2 \|\nabla^3 v\|_2^2 \le 0\); pressure
contributes zero by integration by parts combined with
incompressibility. The new contribution comes from advection.

\subsubsection{5.2 The vortex-stretching
term}\label{the-vortex-stretching-term}

The advective contribution to \(dL/dt\) reduces, after standard
rearrangement using incompressibility, to the canonical
vortex-stretching form:

\[\left(\frac{dL}{dt}\right)_\mathrm{adv} \;\propto\; \int \boldsymbol{\omega} \cdot (S\, \boldsymbol{\omega}) \, dV,\]

where \(\boldsymbol{\omega} = \nabla \times \mathbf{v}\) is the
vorticity and \(S_{ij} = \frac{1}{2}(\partial_i v_j + \partial_j v_i)\)
is the symmetric strain-rate tensor. This is the
\textbf{vortex-stretching term} as it appears in standard
turbulence-analysis texts. (Conventions for the proportionality constant
vary across the literature based on whether \(L\) is defined with or
without a factor of one-half and whether enstrophy or gradient-norm is
used; the physical content of the indefinite-sign vortex-stretching
contribution is convention-independent.)

\subsubsection{5.3 Unique indefinite-sign
contribution}\label{unique-indefinite-sign-contribution}

In the aggregate \(dL/dt\) for full NS + R1,

\[\frac{dL}{dt} = \underbrace{-\nu\|\nabla^2 v\|_2^2}_{\le 0} + \underbrace{-\kappa\mu_{V1}\ell_P^2 \|\nabla^3 v\|_2^2}_{\le 0} + \underbrace{0}_\text{pressure} + \underbrace{\int \boldsymbol{\omega} \cdot S\boldsymbol{\omega} \, dV \cdot (\mathrm{const})}_\text{indefinite},\]

three terms are non-positive (viscous and R1 dissipative; pressure
zero); only the advective vortex-stretching term has indefinite sign.
\textbf{Any potential growth of the gradient norm in 3D NS + R1 must
source from this term alone.} All other contributions to \(dL/dt\) are
dissipative or neutral.

\subsubsection{5.4 2D versus 3D contrast}\label{d-versus-3d-contrast}

The vortex-stretching term
\(\int \boldsymbol{\omega} \cdot S\boldsymbol{\omega} \, dV\) depends on
the alignment of vorticity with strain-rate tensor eigenvectors. The
strain \(S\) has three real eigenvalues
\(\lambda_1 \ge \lambda_2 \ge \lambda_3\) with
\(\lambda_1 + \lambda_2 + \lambda_3 = 0\) by tracelessness
(incompressibility). Vorticity aligned with the
\(\lambda_1\)-eigenvector gives positive integrand locally --- vortex
stretching amplifies vorticity. Vorticity aligned with the
\(\lambda_3\)-eigenvector gives negative integrand locally --- vortex
compression diminishes vorticity. The integrated value is generically
non-zero and indefinite-sign.

In two spatial dimensions, the situation is structurally different.
Vorticity in 2D has only the out-of-plane component
\(\boldsymbol{\omega} = \omega_z \hat{z}\), while the strain-rate tensor
has only in-plane components. Therefore
\(\boldsymbol{\omega} \cdot S\boldsymbol{\omega} = 0\) identically ---
the vortex-stretching term vanishes in 2D.

This is the structural reason 2D NS has Leray-class global smooth
solutions while 3D NS smoothness remains the open Clay problem.
\textbf{The dimensional asymmetry between 2D and 3D in the smoothness
question is not accidental; it is a direct consequence of the
dimension-specific behavior of the advective vortex-stretching content.}
ED's framework makes this asymmetry structurally intelligible: the
obstruction to ED-style monotone gradient-norm decay vanishes in 2D and
is sign-indefinite in 3D.

\subsubsection{5.5 Localization of the
obstruction}\label{localization-of-the-obstruction}

The aggregate finding of Section 5: \textbf{the advective convective
derivative is the unique structural feature of full NS + R1 whose
contribution to the gradient-norm Lyapunov derivative can have positive
sign in 3D.} Pressure and incompressibility do not break the Lyapunov;
they are dissipative-neutral. Viscous diffusion is dissipative. R1's
substrate-scale stabilization is dissipative with sign FORCED positive.
Only advection --- specifically, its vortex-stretching content in 3D ---
produces indefinite-sign contributions.

The structural obstruction to closing Clay-NS via the R1 mechanism alone
is therefore \textbf{localized at the advective term}. This sets up the
architectural decomposition pursued in Section 7.

\begin{center}\rule{0.5\linewidth}{0.5pt}\end{center}

\subsection{6. Three-Angle Convergence on
Advection-is-Non-ED}\label{three-angle-convergence-on-advection-is-non-ed}

\subsubsection{6.1 Architectural convergence
(NS-2.08)}\label{architectural-convergence-ns-2.08}

The first lens identifying advection as structurally non-ED is
architectural. NS-2.08 §5 catalogues the structural features of full NS
that lie outside ED's canonical PDE channels. The catalogue identifies
three fluid-mechanical additions: pressure as Lagrange multiplier, the
advective convective derivative as kinematic coupling between velocity
components, and the incompressibility constraint
\(\nabla \cdot \mathbf{v} = 0\) as a holonomic commitment. Of these
three, advection is the structural feature most strongly outside the
canon: pressure and incompressibility have no architectural P-level
counterparts but contribute zero to dynamical Lyapunov analyses (see
Sections 4--5), while advection's kinematic coupling has no
architectural P-level counterpart \emph{and} contributes the
indefinite-sign vortex-stretching content.

The architectural lens identifies advection as structurally non-ED
because the kinematic-coupling structure of
\((\mathbf{v} \cdot \nabla)\mathbf{v}\) has no analog among the
canonical channels (mobility, penalty, participation). It is
fluid-mechanical-specific structural content, not derivable from any
canonical principle.

\subsubsection{6.2 Dynamical convergence
(NS-Smooth-3)}\label{dynamical-convergence-ns-smooth-3}

The second lens is dynamical. NS-Smoothness Section 5 (NS-Smooth-3 in
the program memo system) computes the gradient-norm Lyapunov derivative
for full NS + R1 explicitly and identifies the advective
vortex-stretching term as the unique indefinite-sign contribution.
Pressure contributes zero; viscous and R1 contributions are manifestly
non-positive; only advection's vortex-stretching can drive growth of the
gradient norm. The Lyapunov analysis localizes the obstruction at the
advective term independently of the architectural-catalogue argument.

The dynamical lens identifies advection as structurally non-ED because
it alone breaks the gradient-norm Lyapunov's monotonicity in 3D --- a
property that ED's R1 mechanism would otherwise enforce.

\subsubsection{6.3 Spectral convergence
(NS-Turb-4)}\label{spectral-convergence-ns-turb-4}

The third lens is spectral. NS-Turb's analysis of the P7 ↔
NS-turbulence-cascade mapping computes the Fourier-space interaction
coefficient of the advective term and finds it to be

\[M_{ijm}(\mathbf{k}) = -i k_j P_{im}(\mathbf{k}), \qquad P_{im}(\mathbf{k}) = \delta_{im} - \frac{k_i k_m}{k^2},\]

with transport-directional structure (via \(k_j\)) and incompressibility
projection (via \(P_{im}\)). This index structure is \emph{asymmetric}
and \emph{transport-directional} --- fundamentally different from ED's
P7 nonlinearity, which is symmetric quadratic-in-gradients
\(M'(\rho) |\nabla\rho|^2\) at the canonical-PDE level. The advective
bilinear-with-projection structure cannot be absorbed into P7-class
symmetric-quadratic Fourier mapping.

The spectral lens identifies advection as structurally non-ED because
its index structure is incompatible with the canonical
nonlinear-coupling form provided by P7.

\subsubsection{6.4 Robustness across analytical
frameworks}\label{robustness-across-analytical-frameworks}

The three lenses operate at three different mathematical levels ---
architectural canon-membership, dynamical Lyapunov-derivative-sign, and
spectral Fourier-mode-coupling --- and use disjoint analytical
machinery. Yet each identifies the same physical feature (advection's
transport-directional, asymmetric, projected index structure) as the
locus of the ED↔NS structural mismatch.

The independence of the three lenses is essential. If only one analysis
identified advection as non-ED, the finding would be susceptible to the
suspicion that the analytical framework itself was unsuited to the
question. Three independent frameworks identifying the same feature
establishes the finding as structural rather than methodological ---
robust across analytical lenses.

This three-angle convergence is the most consistent and load-bearing
structural finding in the ED Navier--Stokes program. It is invoked in
Section 7 as the load-bearing input for the Intermediate Path C verdict.

\begin{center}\rule{0.5\linewidth}{0.5pt}\end{center}

\subsection{7. Intermediate Path C
(Clay-Relevance)}\label{intermediate-path-c-clay-relevance}

\subsubsection{7.1 R1 positive side}\label{r1-positive-side}

ED's canonical architectural content supplies a real Clay-NS-relevant
regularizing mechanism. The R1 mechanism --- the form-FORCED
\(-\kappa\mu_{V1}\ell_P^2\nabla^4 v\) stabilization arising from V1's
finite-width vacuum kernel under multi-scale expansion --- combined with
standard viscous diffusion produces, in the counterfactual ED-only NS
lacking the advective term, a strictly monotonically-decaying
gradient-norm Lyapunov:

\[\frac{dL}{dt}\bigg|_\text{ED-only NS} = -\nu\|\nabla^2 v\|_2^2 - \kappa\mu_{V1}\ell_P^2 \|\nabla^3 v\|_2^2 \;\le\; 0.\]

By standard parabolic-regularity theory, ED-only NS has global smooth
solutions on \(\mathbb{R}^3\) for finite-energy initial data. \emph{If}
3D Navier--Stokes lacked the advective convective derivative, ED's
architectural content would unconditionally close the Clay-NS smoothness
question.

\subsubsection{7.2 Advection negative
side}\label{advection-negative-side}

The actual obstruction to closing Clay-NS in 3D is the advective
convective derivative's vortex-stretching content
\(\int \boldsymbol{\omega} \cdot S\boldsymbol{\omega} \, dV\), which is
the unique indefinite-sign contribution to \(dL/dt\) in full NS + R1.
This term vanishes identically in 2D (where vortex-stretching does not
exist) and is sign-indefinite in 3D.

The advective term is structurally non-ED at three independent
program-level analyses (Section 6): architectural (advection is
fluid-mechanical addition not native to canonical channels), dynamical
(advection alone breaks the gradient-norm Lyapunov via
vortex-stretching), and spectral (advection's bilinear-with-projection
Fourier structure is incompatible with P7-class symmetric-quadratic
mapping). The three-angle convergence makes the advection-as-non-ED
finding robust across analytical lenses.

\subsubsection{7.3 Structural decomposition of smoothness
difficulty}\label{structural-decomposition-of-smoothness-difficulty}

The two sides combine into the \textbf{Intermediate Path C structural
decomposition}:

{\def\LTcaptype{none} % do not increment counter
\begin{longtable}[]{@{}
  >{\raggedright\arraybackslash}p{(\linewidth - 6\tabcolsep) * \real{0.2500}}
  >{\raggedright\arraybackslash}p{(\linewidth - 6\tabcolsep) * \real{0.2500}}
  >{\raggedright\arraybackslash}p{(\linewidth - 6\tabcolsep) * \real{0.2500}}
  >{\raggedright\arraybackslash}p{(\linewidth - 6\tabcolsep) * \real{0.2500}}@{}}
\toprule\noalign{}
\begin{minipage}[b]{\linewidth}\raggedright
Component
\end{minipage} & \begin{minipage}[b]{\linewidth}\raggedright
Source
\end{minipage} & \begin{minipage}[b]{\linewidth}\raggedright
Sign
\end{minipage} & \begin{minipage}[b]{\linewidth}\raggedright
Status
\end{minipage} \\
\midrule\noalign{}
\endhead
\bottomrule\noalign{}
\endlastfoot
Viscous diffusion \(-\nu\|\nabla^2 v\|^2\) & Standard NS & Dissipative
(\(\le 0\)) & Standard \\
R1 stabilization \(-\kappa\mu_{V1}\ell_P^2\|\nabla^3 v\|^2\) & ED
canonical & Dissipative (\(\le 0\)); sign FORCED & Form-FORCED ED
architecture \\
Pressure & Lagrange multiplier & Zero & Fluid-mechanical addition \\
Advection vortex-stretching & Fluid-mechanical addition &
Indefinite-sign & Non-ED (three-angle convergence) \\
\end{longtable}
}

The decomposition makes the Clay-NS difficulty intelligible. ED supplies
real \emph{regularizing infrastructure} (R1), form-FORCED,
sign-FORCED-positive, dissipative. The structural feature
\emph{breaking} gradient-norm monotonicity in 3D is \emph{not} in ED's
canonical content; it is the fluid-mechanical-specific advective
coupling. The \emph{quantitative competition} between R1's dissipative
content (dominant only at substrate scales \(\sim \ell_P\)) and
advective vortex-stretching (active at intermediate scales between flow
scale and substrate) determines whether smoothness preserves or breaks.
\textbf{This competition is INHERITED on both sides} --- depends on V1's
specific G-function profile (inherited per arc-N N.4) and on standard
kinetic-theory super-Burnett magnitude (inherited via material kinetic
parameters). Neither magnitude is canonically fixed; therefore neither
is the dominant in any specific blow-up scenario.

\subsubsection{7.4 What ED explains}\label{what-ed-explains}

The structural decomposition explains:

\begin{itemize}
\tightlist
\item
  \textbf{Why the Navier--Stokes form appears.} NS reproduces from
  substrate primitives via two concordant routes; the viscous content is
  canonical ED architecture.
\item
  \textbf{Why 2D NS is globally smooth.} Vortex-stretching vanishes
  identically in 2D; the gradient-norm Lyapunov is monotone-decreasing
  exactly as in ED-only NS.
\item
  \textbf{Why 3D NS smoothness is structurally hard.} The obstructing
  structural feature lies outside ED's canonical regularizing
  architecture; the canon does not natively absorb advection.
\item
  \textbf{Where the obstruction is localized.} Uniquely at the advective
  convective derivative --- confirmed at three independent analytical
  levels.
\item
  \textbf{Why R1's substrate-scale stabilization is real but not
  unconditionally sufficient.} R1 is form-FORCED and
  sign-FORCED-positive, but its coefficient is INHERITED at the value
  level and the term is suppressed at intermediate scales where
  advective vortex-stretching is most active.
\end{itemize}

\subsubsection{7.5 What ED does not
explain}\label{what-ed-does-not-explain}

The structural decomposition does not deliver:

\begin{itemize}
\tightlist
\item
  \textbf{Whether 3D NS solutions blow up at finite time or remain
  smooth globally.} The Clay-NS open question remains open; ED does not
  resolve which side wins quantitatively.
\item
  \textbf{Numerical critical Reynolds numbers} for any specific
  transition (laminar-to-turbulent, pipe flow Re\_c ≈ 2300,
  boundary-layer transition, etc.).
\item
  \textbf{Specific blow-up criteria} beyond the Beale--Kato--Majda
  framework.
\item
  \textbf{The detailed cascade structure} of developed turbulence
  (Kolmogorov 5/3 spectrum, intermittency exponents, anomalous scaling).
\end{itemize}

ED's contribution to the Clay-NS question is
\textbf{structural-decompositional, not quantitative-resolutional}. This
is the substantive Intermediate Path C content. The framework neither
solves Clay-NS nor is irrelevant to it; it provides a partial structural
framework that explains why the difficulty exists where it does without
resolving how the difficulty plays out quantitatively.

\begin{center}\rule{0.5\linewidth}{0.5pt}\end{center}

\subsection{8. Turbulence}\label{turbulence}

\subsubsection{8.1 H1 trivial}\label{h1-trivial}

The first hypothesis tested in the NS-Turb arc is H1: ED's P7
nonlinear-triad-coupling structure shares generic triadic Fourier
structure with the NS turbulent cascade. This succeeds trivially. Both
P7's harmonic generation \(M'(\rho)|\nabla\rho|^2\) and NS's advective
term \((\mathbf{v} \cdot \nabla)\mathbf{v}\) produce triadic Fourier
interactions \(\mathbf{k}_1 + \mathbf{k}_2 + \mathbf{k}_3 = 0\) ---
momentum conservation in mode-mode interactions. This is generic to any
quadratic-or-higher-in-field nonlinearity in Fourier space; it carries
no ED-specific architectural insight beyond ``NS has a quadratic
nonlinearity.''

H1 closure: succeeds trivially; provides no ED-specific structural
content beyond the generic Fourier-triadic structure.

\subsubsection{8.2 H2 partial}\label{h2-partial}

The second hypothesis tested is H2: ED's canonical 3--6\% P7 amplitude
ratio (Universal Invariants §3) corresponds to a specific
turbulence-cascade observable. This succeeds \emph{partially} in a
restricted regime.

In the inertial range of developed turbulence, dimensionless
triad-transfer observables (fractional transfer per eddy turnover,
normalized triad correlation coefficient) have values
\(\tilde f \sim 0.1\)--\(0.3\) for active cascade triads --- 2 to 10
times larger than P7's 3--6\%. H2 fails in the cascade regime.

In the \emph{forced-response} regime of viscous-laminar
single-wavenumber-forced flow, by contrast, weakly-nonlinear analysis
gives a third-harmonic amplitude ratio

\[R_\mathrm{forced} = \frac{|u(3k_0)|}{|u(k_0)|} \sim \varepsilon^2,\]

with \(\varepsilon = u(k_0)/(\nu k_0)\) a dimensionless forcing-mode
parameter. For moderate weakly-nonlinear
\(\varepsilon \sim 0.15\)--\(0.3\), \(R_\mathrm{forced}\) falls in the
range 0.02--0.09 --- overlapping P7's 3--6\% range cleanly.

H2 holds in this restricted forced-response regime as a substantive but
limited correspondence. The match is at order-of-magnitude /
window-overlap level rather than at precise-prefactor level, and it does
not extend to the developed-cascade regime that defines turbulence
empirically. The match is \textbf{coincidental rather than structural}
--- both phenomena have triadic Fourier structure with weak coupling,
producing similar amplitude scaling in the appropriate regimes, but the
underlying dynamics differ in mechanism (driven vs.~cascade), polynomial
order (cubic vs.~bilinear), and index structure (symmetric-gradient
vs.~transport-directional).

\subsubsection{8.3 H3 fail}\label{h3-fail}

The third hypothesis tested is H3: ED's P7 supplies an architectural
template for the developed-cascade turbulence structure of 3D NS. This
fails. Three structural mismatches prevent P7 from templating the
cascade:

\begin{enumerate}
\def\labelenumi{\arabic{enumi}.}
\item
  \textbf{Mechanism mismatch.} P7 is forced harmonic generation by
  external driving on a fundamental; turbulence cascade is free
  conservative redistribution across modes. The triadic conservation
  identity
  \(S(\mathbf{k}|\mathbf{p},\mathbf{q}) + S(\mathbf{p}|\mathbf{q},\mathbf{k}) + S(\mathbf{q}|\mathbf{k},\mathbf{p}) = 0\)
  holds in the inertial range, distinguishing free-cascade physics from
  driven harmonic-generation physics.
\item
  \textbf{Polynomial-order mismatch.} P7 nonlinearity is cubic in field
  (third harmonic at first nonlinear order); NS advection is bilinear in
  field (third harmonic at second nonlinear order). Same numerical
  magnitude requires different effective forcing intensities.
\item
  \textbf{Index-structure asymmetry.} P7 is symmetric
  quadratic-in-gradients; NS advection is bilinear field-times-gradient
  with directional transport and incompressibility projection. The
  asymmetry persists at all amplitudes.
\end{enumerate}

H3 fails. ED's P7 does not architecturally template developed-cascade
turbulence.

\subsubsection{8.4 Spectral incompatibility with
P7}\label{spectral-incompatibility-with-p7}

The deepest finding from NS-Turb is the spectral incompatibility of the
advective interaction coefficient with P7-class structure. Computing the
NS triadic interaction coefficient explicitly,

\[M_{ijm}(\mathbf{k}) = -i k_j P_{im}(\mathbf{k}),\]

with transverse projector
\(P_{im}(\mathbf{k}) = \delta_{im} - k_i k_m/k^2\) enforcing
incompressibility, reveals a transport-directional, asymmetric,
projected structure. P7's nonlinearity in Fourier space, by contrast,
produces a symmetric-in-permutation interaction coefficient appropriate
to gradient self-coupling. The two are structurally different; no
constitutive choice for \(M(\rho)\) in P7 produces the
transport-directional content of NS advection.

The spectral incompatibility is the third angle in the three-angle
convergence on advection-as-non-ED (Section 6).

\subsubsection{8.5 Why ED has no turbulence-cascade
template}\label{why-ed-has-no-turbulence-cascade-template}

The aggregate verdict of NS-Turb is that \textbf{ED's canonical content
does not architecturally template developed-cascade turbulence}. The
cascade dynamics --- the substantive content of ``what makes turbulence
turbulence'' --- is built on the advective vortex-stretching mechanism,
which is non-ED at architectural, dynamical, and spectral levels. ED's
contribution to fluid mechanics in the turbulence regime is therefore
limited:

\begin{itemize}
\tightlist
\item
  \textbf{Forced-response harmonic generation} in viscous-laminar
  regime: ED-architectural restricted-regime prediction (H2 partial).
\item
  \textbf{Cascade structure} (Kolmogorov spectrum, intermittency,
  anomalous scaling): no ED-architectural prediction (H3 fail).
\item
  \textbf{Smoothness / blow-up} (Clay-NS): Intermediate Path C;
  obstruction is the same advective vortex-stretching that NS-Turb
  identifies as breaking P7-class structure.
\end{itemize}

The three program-level findings --- Architectural Canon Vector
Extension's catalogue of advection as non-ED, NS-Smoothness's
localization of the gradient-norm Lyapunov obstruction, and NS-Turb's
spectral incompatibility --- are mutually reinforcing. The same
structural feature is the locus of ED's bounded reach into
Navier--Stokes phenomena.

\begin{center}\rule{0.5\linewidth}{0.5pt}\end{center}

\subsection{9. Synthesis}\label{synthesis}

\subsubsection{9.1 Unified structural
picture}\label{unified-structural-picture}

The five closed Navier--Stokes arcs assemble into a coherent unified
structural picture. Event Density's relationship to classical fluid
mechanics is neither one of identity nor of irrelevance: it is partial
and structurally honest.

\begin{itemize}
\tightlist
\item
  \textbf{Form derivation:} the Navier--Stokes form follows from
  substrate primitives via two concordant routes (NS-2). Both routes
  produce the same Newtonian-fluid form.
\item
  \textbf{Dimensional forcing:} the spatial dimension D = 3+1 is forced
  by ED's framework when combined with observed gravity (NS-1, Path
  B-strong).
\item
  \textbf{Architectural classification:} NS is partially
  ED-architectural --- its viscous content is canonical ED
  vector-extended; its pressure, advection, and incompressibility
  content is fluid-mechanical-specific structural addition (NS-2.08
  catalogue).
\item
  \textbf{Clay-relevance:} ED contains a real, form-FORCED regularizing
  mechanism (R1) that closes ED-only NS smoothness unconditionally and
  that is suppressed but structurally present in full NS (NS-3 +
  NS-Smoothness).
\item
  \textbf{Turbulence cascade:} no architectural template; the cascade
  mechanism is non-ED (NS-Turb).
\item
  \textbf{Three-angle convergence:} advection is structurally non-ED at
  architectural, dynamical, and spectral levels independently --- the
  most robust structural finding in the program.
\end{itemize}

\subsubsection{9.2 ED's position relative to classical
Navier--Stokes}\label{eds-position-relative-to-classical-navierstokes}

The unified picture positions ED relative to classical NS as follows.
Classical NS is a phenomenological framework: it postulates the equation
form, takes the dimensional structure as observational, and treats the
smoothness question as an open mathematical problem. ED is an
architectural framework: it derives the equation form (partly), forces
the dimensional structure, and decomposes the smoothness question into a
part it explains (the regularizing infrastructure) and a part it does
not (the advective obstruction). The two frameworks agree on
observational content --- ED reproduces standard Newtonian-fluid NS at
NS scales --- and differ on structural origin.

Classical NS is the empirically successful phenomenology. ED is the
architectural account underneath that phenomenology, exposing which
structural features are forced by substrate-level commitments and which
are fluid-mechanical-specific additions inherited from the kinematic
structure of incompressible flow.

\subsubsection{9.3 ED's architectural reach and
limits}\label{eds-architectural-reach-and-limits}

The reach of ED's architectural content into Navier--Stokes phenomena is
substantial but bounded. The canonical canon delivers the viscous
content of NS, the dimensional structure D = 3+1, the substrate-cutoff
regularization R1, the gradient-norm Lyapunov for ED-only NS, the
2D-versus-3D smoothness asymmetry, and the structural localization of
the Clay obstruction. These are all forced theorems of the architectural
canon plus its vector-extension.

The canon does not deliver Clay-NS resolution, critical Reynolds
numbers, turbulence cascade architecture, or specific blow-up criteria.
These lie outside the architectural reach. The boundary of ED's
architectural content is itself structurally identified: it is the
boundary of what the canonical principles P1--P7 plus the partial
vector-extension framework can absorb. Advection lies outside this
boundary at three independent levels.

This honest delineation of reach and limits is the substantive
structural content of the synthesis: ED is a partial architectural
account of NS, with the boundary of its reach explicitly identified and
characterized.

\subsubsection{9.4 Relationship to the P4-NN rheology
paper}\label{relationship-to-the-p4-nn-rheology-paper}

This synthesis paper covers the foundational structural content of ED's
relationship to Navier--Stokes --- form derivation, dimensional forcing,
Clay-relevance decomposition, turbulence verdict. The companion P4-NN
paper (\emph{ED Mobility Saturation Predicts Non-Newtonian Rheology})
covers the empirical-applications content of ED's reach into
non-Newtonian fluid mechanics: Krieger--Dougherty divergence,
discontinuous shear-thickening, Cross/Carreau shear-thinning,
Maxwell-class viscoelasticity, and the canonical \(\beta = 2.0\)
exponent prediction supported within \(1\sigma\) by the Universal
Mobility Law's ten-system empirical anchor.

The two papers together cover the program's published-grade content on
classical fluids: foundational structural questions in this paper,
empirical applications in P4-NN. They share architectural content (the
same canonical canon, the same form-FORCED / value-INHERITED
methodology) but address different levels of the fluid-mechanics
hierarchy. Both are grounded in the same Universal Mobility Law
\(D(c) = D_0(1 - c/c_\mathrm{max})^\beta\) --- the same architectural
principle (P4 mobility capacity bound) that produces non-Newtonian
rheology classes empirically also produces, in the NS context, the
dissipative content of the Newtonian viscous term.

The two papers complement, not replace, each other.

\begin{center}\rule{0.5\linewidth}{0.5pt}\end{center}

\subsection{10. Conclusion}\label{conclusion}

\subsubsection{10.1 Summary}\label{summary}

We have presented a unified architectural account of the Navier--Stokes
equations within the Event Density program. From substrate-level
commitments, the canonical PDE's vector-extension produces the viscous
content of standard Newtonian-fluid NS, with the dimensional structure D
= 3+1 forced via a Path B-strong argument combining architectural and
empirical-consistency lines. Substrate-level coarse-graining of chain
dynamics produces the same NS form via a methodologically distinct
route, and the two derivations are concordant. Pressure, the advective
convective derivative, and the incompressibility constraint are
catalogued explicitly as fluid-mechanical-specific structural additions
not native to ED canonical channels.

The architectural canon supplies a real, form-FORCED regularizing
mechanism --- the R1 substrate-scale stabilization arising from V1's
finite-width vacuum kernel --- that produces, in the counterfactual
ED-only NS lacking advection, strictly monotonically-decaying
gradient-norm Lyapunov decay and global smooth solutions on
\(\mathbb{R}^3\). The advective vortex-stretching content of full 3D NS
is the unique structural feature whose contribution to the gradient-norm
Lyapunov derivative can be sign-indefinite; it vanishes identically in
2D, accounting for the dimensional asymmetry of NS smoothness. Three
independent program-level analyses (architectural, dynamical, spectral)
converge on advection as structurally non-ED --- the most robust
structural finding in the program.

The Intermediate Path C verdict follows: ED neither solves the Clay-NS
smoothness problem nor is irrelevant to it. ED supplies a partial
structural framework that explains why the difficulty exists where it
does, decomposing the Clay-NS question into ED-canonical regularizing
infrastructure plus non-ED structural obstruction. The decomposition
explains the 2D-versus-3D smoothness asymmetry, localizes the 3D
obstruction at the advective term, and is robust across architectural,
dynamical, and spectral lenses. It does not resolve which side wins
quantitatively; that competition is INHERITED on both sides.

The closure of NS-Turb at H1 trivial / H2 partial / H3 fail establishes
that ED supplies no architectural template for developed-cascade
turbulence. The cascade mechanism is built on the advective
vortex-stretching content that the three-angle convergence identifies as
non-ED. The reach of ED's architectural content into Navier--Stokes
phenomena is therefore substantial but bounded: it covers form
derivation, dimensional forcing, structural smoothness decomposition,
and the 2D-versus-3D asymmetry; it does not cover Clay-NS resolution,
critical Reynolds numbers, or turbulence cascade architecture.

\subsubsection{10.2 Future directions}\label{future-directions}

Three concrete directions emerge as candidate future arcs:

\begin{itemize}
\tightlist
\item
  \textbf{Path A B2 promotion via diffusion-coarse-graining theorem.}
  NS-1.03's identification gap (chain-dynamics-as-diffusion at
  coarse-grained scale) would close if a primitive-level
  coarse-graining-to-diffusion theorem for ED chain dynamics were
  established. This would promote NS-1's architectural-suggestive d ≤ 3
  to a strict architectural theorem.
\item
  \textbf{Path C+ promotion via R1-versus-Burnett quantitative
  dominance.} NS-3's Intermediate Path C verdict could promote toward
  Path C+ (clean Clay-NS resolution) if a structural argument
  established R1's substrate-scale stabilization as quantitatively
  dominant over standard kinetic-theory super-Burnett destabilization
  through scales. This is value-level work depending on parameters
  INHERITED on both sides.
\item
  \textbf{Tensor-extension of the Architectural Canon.} Rheology classes
  that are non-canonical at the current canon (Oldroyd-B, Giesekus,
  FENE-P) require conformation-tensor structure beyond the canonical
  scalar PDE. A formal tensor-extension parallel to the vector-extension
  framework would absorb these into ED's architectural reach.
\end{itemize}

Each is substantive future work, parallel in scope to the closed arcs of
the present paper. None is required for the Intermediate Path C verdict.

\subsubsection{10.3 Program-level
implications}\label{program-level-implications}

This paper closes the structural-foundations content of the Event
Density program's relationship to Navier--Stokes. With the closure of
NS-1, NS-2, NS-3, NS-Smoothness, and NS-Turb, the program has a complete
picture of how ED relates to classical fluid mechanics, the Clay-NS
smoothness problem, and the structural status of turbulence. Combined
with the empirical-applications content of the P4-NN companion paper,
the program's published-grade fluid-mechanics content is now
consolidated.

The framework's posture is the same as in the broader ED program:
substrate primitives plus structural canon principles produce, as forced
theorems, the foundational structures of nominally separate domains of
physics. For Navier--Stokes the result is partial --- the canon supplies
the viscous content and the dimensional forcing; advection is identified
as structurally outside the canon. The boundary of ED's reach is itself
structurally identified, not hand-waved. That honest delineation of
reach and limits is the substantive contribution of the synthesis.

Whether the structural decomposition presented here becomes part of a
future Clay-NS resolution, or whether the Clay-NS question is resolved
by methods entirely outside ED's architectural framework, is an
empirical-mathematical question that future analytical work must answer.
ED's case is on the table.

\begin{center}\rule{0.5\linewidth}{0.5pt}\end{center}

\subsection{Appendix A --- Key
Equations}\label{appendix-a-key-equations}

For reference, the load-bearing equations of the synthesis are collected
here.

\textbf{Canonical Event Density PDE:}

\[\partial_t \rho = D \cdot F[\rho] + H \cdot v, \qquad F[\rho] = M(\rho)\nabla^2\rho + M'(\rho)|\nabla\rho|^2 - P(\rho).\]

\textbf{ED-only NS equation (counterfactual; without advection):}

\[\rho \, \partial_t v_i = \mu \nabla^2 v_i - \kappa \mu_{V1} \ell_P^2 \nabla^4 v_i - \partial_i p, \qquad \nabla \cdot \mathbf{v} = 0.\]

\textbf{Full Navier--Stokes with R1 stabilization:}

\[\rho \, \partial_t v_i + \rho \, (v \cdot \nabla) v_i = \mu \nabla^2 v_i - \kappa \mu_{V1} \ell_P^2 \nabla^4 v_i - \partial_i p, \qquad \nabla \cdot \mathbf{v} = 0.\]

\textbf{Gradient-norm Lyapunov:}

\[L(t) = \frac{1}{2} \|\nabla \mathbf{v}(t)\|_2^2.\]

\textbf{Gradient-norm Lyapunov derivative for ED-only NS:}

\[\frac{dL}{dt}\bigg|_\text{ED-only NS} = -\nu \int |\nabla^2 v|^2 \, dV - \kappa \mu_{V1} \ell_P^2 \int |\nabla^3 v|^2 \, dV \le 0.\]

\textbf{Gradient-norm Lyapunov derivative for full NS + R1:}

\[\frac{dL}{dt} = -\nu\|\nabla^2 v\|_2^2 - \kappa\mu_{V1}\ell_P^2 \|\nabla^3 v\|_2^2 + 0 + \int \boldsymbol{\omega} \cdot S\boldsymbol{\omega} \, dV \cdot (\mathrm{const}).\]

\textbf{Vortex-stretching contribution (advective):}

\[\left(\frac{dL}{dt}\right)_\mathrm{adv} \propto \int \boldsymbol{\omega} \cdot (S\,\boldsymbol{\omega}) \, dV.\]

\textbf{NS triadic interaction coefficient (Fourier space):}

\[M_{ijm}(\mathbf{k}) = -i k_j P_{im}(\mathbf{k}), \qquad P_{im}(\mathbf{k}) = \delta_{im} - \frac{k_i k_m}{k^2}.\]

\textbf{Forced-response third-harmonic ratio (NS-Turb-3):}

\[R_\mathrm{forced} = \frac{|u(3k_0)|}{|u(k_0)|} \sim \varepsilon^2, \qquad \varepsilon = \frac{u(k_0)}{\nu k_0}.\]

\begin{center}\rule{0.5\linewidth}{0.5pt}\end{center}

\subsection{Appendix B --- Summary of Closed Navier--Stokes
Arcs}\label{appendix-b-summary-of-closed-navierstokes-arcs}

{\def\LTcaptype{none} % do not increment counter
\begin{longtable}[]{@{}
  >{\raggedright\arraybackslash}p{(\linewidth - 4\tabcolsep) * \real{0.3333}}
  >{\raggedright\arraybackslash}p{(\linewidth - 4\tabcolsep) * \real{0.3333}}
  >{\raggedright\arraybackslash}p{(\linewidth - 4\tabcolsep) * \real{0.3333}}@{}}
\toprule\noalign{}
\begin{minipage}[b]{\linewidth}\raggedright
Arc
\end{minipage} & \begin{minipage}[b]{\linewidth}\raggedright
Headline
\end{minipage} & \begin{minipage}[b]{\linewidth}\raggedright
Status
\end{minipage} \\
\midrule\noalign{}
\endhead
\bottomrule\noalign{}
\endlastfoot
\textbf{NS-1} & D = 3+1 forced via Path B-strong (architectural d ≥ 3 +
architectural-suggestive d ≤ 3 + empirical-consistency d ≤ 3) &
Closed \\
\textbf{NS-2} & Standard Newtonian-fluid NS form derived via two
concordant routes (chain-substrate + ED-PDE-direct); advection /
pressure / incompressibility flagged as fluid-mechanical additions &
Closed \\
\textbf{NS-3} & Intermediate Path C: ED contains real Clay-relevant R1
mechanism (form-FORCED \(-\kappa\mu_{V1}\ell_P^2\nabla^4 v\));
quantitative competition vs.~super-Burnett INHERITED & Closed \\
\textbf{NS-Smoothness} & Clay-relevance decomposition: R1 +
advection-obstruction; ED-only NS has \(dL/dt \le 0\); advective
vortex-stretching is unique indefinite-sign contribution in full NS +
R1; three-angle convergence on advection-is-non-ED & Closed \\
\textbf{NS-Turb} & H1 trivial / H2 partial (forced-response regime,
\(\varepsilon \sim 0.15\)--\(0.3\)) / H3 fail; no architectural template
for developed-cascade turbulence & Closed \\
\end{longtable}
}

\begin{center}\rule{0.5\linewidth}{0.5pt}\end{center}

\subsection{Appendix C --- Architectural Canon
References}\label{appendix-c-architectural-canon-references}

\begin{itemize}
\tightlist
\item
  \textbf{P1 (operator structure):}
  \(F[\rho] = M(\rho)\nabla^2\rho + M'(\rho)|\nabla\rho|^2 - P(\rho)\).
\item
  \textbf{P2 (channel complementarity):}
  \(\partial_t \rho = D \cdot F[\rho] + H \cdot v\), \(D + H = 1\).
\item
  \textbf{P3 (penalty equilibrium):} \(P(\rho^*) = 0\),
  \(P'(\rho^*) > 0\), monostable.
\item
  \textbf{P4 (mobility capacity bound):} \(M(\rho_\mathrm{max}) = 0\),
  \(M(\rho) > 0\) for \(\rho < \rho_\mathrm{max}\).
\item
  \textbf{P5 (participation feedback):}
  \(\dot v = \tau^{-1}(\bar F(\rho) - \zeta v)\).
\item
  \textbf{P6 (damping discriminant):}
  \(D_\mathrm{crit}(\zeta) \approx 0.896\) at \(\zeta = 1/4\).
\item
  \textbf{P7 (nonlinear triad coupling):} harmonic generation at
  canonical 3--6\% from \(M'(\rho)|\nabla\rho|^2\).
\end{itemize}

The Architectural Canon Vector Extension establishes that P1--P7 are
field-type-agnostic at architectural level; vector and tensor field PDEs
satisfying P1--P7 component-wise are architecturally ED. C5 (single
scalar field) binds at the concrete-PDE-level canonical exemplar but not
at the P-level architectural canon. The three-tier classification
distinguishes canonical / fully ED-architectural / partially
ED-architectural PDEs.

\begin{center}\rule{0.5\linewidth}{0.5pt}\end{center}

\subsection{References}\label{references}

\begin{itemize}
\tightlist
\item
  \emph{Event Density: One Substrate, Three Domains --- A Structural
  Account of Quantum Mechanics, Galactic Gravity, and Universal
  Soft-Matter Mobility} (Proxmire 2026). The unified ED program-overview
  paper.
\item
  \emph{Universal Degenerate-Mobility Scaling in Crowded Soft Matter}
  (Proxmire 2026). The Universal Mobility Law empirical-anchor paper.
\item
  \emph{ED Mobility Saturation Predicts Non-Newtonian Rheology --- A
  Fluid-Mechanical Extension of the Universal Mobility Law} (Proxmire
  2026). The companion non-Newtonian rheology paper (P4-NN).
\item
  \emph{The Primitive-Level Arrow of Time in Event Density: The Closure
  of Arc B and Theorem 18} (Proxmire 2026). The V1 retarded-kernel
  structural origin.
\item
  \emph{Structural Foundations of ED-Substrate Gravity: Newton, the
  Transition Scale, the Combination Rule, and the Baryonic Tully--Fisher
  Relation} (Proxmire 2026). The substrate-gravity foundations paper,
  including T19 (\(G = c^3 \ell_P^2/\hbar\)) and T20
  (\(a_0 = c \cdot H_0/(2\pi)\)).
\item
  \emph{A UV-Finite Quantum Field Theory from Event-Density Primitives}
  (Proxmire 2026). Arc Q.8 vacuum factorisation.
\item
  \emph{Gauge Fields as Forced Rule-Type Structure: Theorem 17}
  (Proxmire 2026). Arc Q closure.
\item
  Beale, J.T., Kato, T., Majda, A. (1984). \emph{Remarks on the
  breakdown of smooth solutions for the 3-D Euler equations.} Comm.
  Math. Phys. 94, 61--66.
\item
  Constantin, P., Foias, C. (1988). \emph{Navier-Stokes Equations.}
  University of Chicago Press.
\item
  Frisch, U. (1995). \emph{Turbulence: The Legacy of A. N. Kolmogorov.}
  Cambridge University Press.
\item
  Lions, J.-L. (1969). \emph{Quelques méthodes de résolution des
  problèmes aux limites non linéaires.} Dunod, Paris.
\item
  Pope, S. B. (2000). \emph{Turbulent Flows.} Cambridge University
  Press.
\end{itemize}

\begin{center}\rule{0.5\linewidth}{0.5pt}\end{center}

\emph{Event Density is an open research framework being developed
publicly by independent physicist Allen Proxmire. The full technical
program --- primitives, derivations, simulations, twenty-one forced
theorems plus the ED Combination Rule, and pre-registered experimental
tests across three domains --- is available at
github.com/allen-proxmire, with archival DOIs through Zenodo. The
complete Navier--Stokes program memos (NS-1 through NS-Turb-5) are in
the program repository under theory/Navier Stokes/. April 2026.}

\begin{center}\rule{0.5\linewidth}{0.5pt}\end{center}

\section{Appendix C --- MHD Extension and ED-I-06 Ontological
Integration}\label{appendix-c-mhd-extension-and-ed-i-06-ontological-integration}

\emph{Companion appendix to} The Architectural Foundations of
Navier-Stokes in Event Density: Form Derivation, Clay-Relevance
Decomposition, and the Structural Status of Turbulence.

\begin{center}\rule{0.5\linewidth}{0.5pt}\end{center}

\subsection{C.1 Purpose}\label{c.1-purpose}

This appendix extends the architectural account of the main paper in two
complementary directions.

First, it extends the NS architectural classification (Sections §3--§6
of the main paper) to \textbf{incompressible magnetohydrodynamics
(MHD)}, the canonical theory of electrically conducting fluids that
couples Navier-Stokes to Maxwell's equations through Ohm's law and the
induction equation. The appendix shows that the canonical /
non-canonical boundary established for pure NS generalizes cleanly to
MHD, with two structural refinements: every electromagnetic-side
addition introduced by the MHD upgrade is ED-canonical, and the new
fluid-mechanical-specific structural commitments (induction-equation
kinematic term, Ohm's-law kinematic component) belong to a single
transport-kinematic class shared with NS advection.

Second, it integrates the \textbf{ED-I-06 \emph{Fields and Forces}
ontology} (Proxmire, \emph{Fields and Forces in Event Density}, 2026) as
the conceptual framework that explains \emph{why} the canonical /
non-canonical boundary lives where it does. The ED-I-06 ontology
supplies a forces-vs-frame-kinematic reading under which the structural
classifications of NS-2.08, NS-Smoothness, NS-Turbulence, and the MHD
arc all read as instances of a single ontological distinction.

The appendix delivers a final architectural classification table
covering all eleven content items of the standard incompressible MHD
system, with both structural and ontological status indicated for each.
The main results of the NS Synthesis Paper remain unchanged; the
ontological integration provides ED-I-06-grounded framing for those
results without revising any technical conclusion.

\begin{center}\rule{0.5\linewidth}{0.5pt}\end{center}

\subsection{C.2 ED-I-06 Ontology: Forces, Fields, and Participation
Structures}\label{c.2-ed-i-06-ontology-forces-fields-and-participation-structures}

ED-I-06 reframes fields and forces in terms of stable participation
structures. Three structural classes of field-like participation arise
naturally from the internal organization of participation channels:

\begin{itemize}
\tightlist
\item
  \textbf{Directional fields} --- channels with a preferred orientation.
  ED-I-06 cites magnetic structure, spin textures, and
  \emph{vorticity-like fluid structures} explicitly (§3). In the NS /
  MHD context, the velocity field \(\mathbf{v}\), the vorticity
  \(\boldsymbol{\omega}\), the magnetic field \(\mathbf{B}\), and the
  gauge field \(A_\mu\) (formalized as \(\tau_g\)'s participation
  measure under T17) are all directional fields.
\item
  \textbf{Scalar fields} --- gradients in participation density.
  Pressure \(p\) and mass density \(\rho\) in fluid mechanics, electric
  and chemical potentials, energy landscapes --- all are emergent
  properties of how participation density varies across space (ED-I-06
  §4).
\item
  \textbf{Curvature-like fields} --- pre-geometric participation
  patterns where propagation paths are biased in geometric-bending-like
  ways. ED-I-06 §5 designates this class as the conceptual bridge to
  spacetime emergence in ED-10. Substrate-gravity content (Theorems 19,
  20, ED Combination Rule, T21) belongs to this class at substrate
  level.
\end{itemize}

\textbf{Forces.} ED-I-06 defines forces as \textbf{biases in
participation flow sourced by stable participation structures} (§6). The
underlying principle: \emph{micro-events follow the most stable
participation channels available to them.} When directional, scalar, or
curvature-like structures bias propagation, the system experiences what
is conventionally interpreted as a force.

\textbf{Non-forces.} Two structurally distinct classes of non-force
content can appear in continuum equations:

\begin{itemize}
\tightlist
\item
  \textbf{Transport-kinematic frame terms} --- bilinear couplings
  between transport fields that arise specifically when dynamics are
  written in a fluid (Eulerian) coordinate system. These are
  coordinate-frame artifacts of the fluid description, not biases
  sourced by stability.
\item
  \textbf{Continuum-imposed constraints} --- fluid-mechanical structural
  commitments (incompressibility, pressure as Lagrange multiplier)
  imposed at continuum level rather than sourced by participation
  structure.
\end{itemize}

The key ontological boundary is:

\begin{quote}
\emph{Canonical ED content corresponds to forces from participation
structures.}

\emph{Non-ED content corresponds to transport-kinematic frame terms and
continuum constraints.}
\end{quote}

This single boundary unifies six previously separate findings of the NS
/ MHD program:

\begin{itemize}
\tightlist
\item
  \textbf{NS-2.08} classified advection, pressure, and incompressibility
  as fluid-mechanical-additions not native to ED canonical channels.
\item
  \textbf{NS-Smoothness} (Section §5 of the main paper) identified the
  advective vortex-stretching
  \(\int\boldsymbol{\omega}\cdot S\boldsymbol{\omega}\,dV\) as the
  unique indefinite-sign contribution to the gradient-norm Lyapunov in
  3D NS.
\item
  \textbf{NS-Turbulence} (Section §6 of the main paper) found the NS
  advective triad to be spectrally incompatible with P7's symmetric
  quadratic class.
\item
  \textbf{NS-MHD-2} identified the Lorentz force as canonical ED via T17
  minimal coupling, with the kinematic \(\mathbf{v}\times\mathbf{B}\)
  component derived from the minimal-coupling form rather than
  separately committed.
\item
  \textbf{NS-MHD-3} established three-angle convergence on the
  induction-equation kinematic term as non-ED at architectural /
  dynamical / spectral levels --- the magnetic-side analogue of NS
  advection.
\item
  \textbf{NS-MHD-4} consolidated the eleven-item architectural
  classification of incompressible MHD.
\end{itemize}

Under the ED-I-06 ontology, all six findings read as instances of the
single forces-vs-frame-kinematic distinction. The transport-kinematic
class (advection, induction kinematic, Ohm kinematic) consists of frame
artifacts of the fluid coordinate system; the canonical-ED class
(viscous diffusion, magnetic diffusion, Lorentz force, Maxwell
structure, R1 substrate cutoff) consists of biases sourced by stable
participation structures (directional fields, the V1 vacuum kernel, the
gauge directional-field \(\tau_g\)).

\begin{center}\rule{0.5\linewidth}{0.5pt}\end{center}

\subsection{C.3 The Incompressible MHD
System}\label{c.3-the-incompressible-mhd-system}

The standard incompressible resistive MHD equations couple Navier-Stokes
to Maxwell's equations via the Lorentz force in the momentum equation
and the induction equation for the magnetic field:

\textbf{Momentum equation:}
\[\rho\bigl(\partial_t\mathbf{v} + (\mathbf{v}\cdot\nabla)\mathbf{v}\bigr) = -\nabla p + \mu\nabla^2\mathbf{v} + (\nabla\times\mathbf{B})\times\mathbf{B} + \rho\mathbf{f}^\mathrm{ext}.\]

\textbf{Induction equation:}
\[\partial_t\mathbf{B} = \nabla\times(\mathbf{v}\times\mathbf{B}) + \eta\nabla^2\mathbf{B}.\]

\textbf{Constraints:}
\[\nabla\cdot\mathbf{v} = 0, \qquad \nabla\cdot\mathbf{B} = 0.\]

The Lorentz force per unit volume in the MHD limit is
\(\mathbf{j}\times\mathbf{B} = \mu_0^{-1}(\nabla\times\mathbf{B})\times\mathbf{B}\),
with quasi-neutrality suppressing the \(\rho_q\mathbf{E}\) contribution
at typical MHD scales. The magnetic diffusivity is
\(\eta = 1/(\sigma\mu_0)\), with \(\sigma\) the electrical conductivity.
The \textbf{ED-augmented momentum equation} carries an additional R1
term \(-\kappa\mu_\mathrm{V1}\ell_P^2\nabla^4\mathbf{v}\) inherited from
the velocity sector (§4 of the main paper, NS-3.01 / NS-Smoothness-2):

\[\rho\bigl(\partial_t\mathbf{v} + (\mathbf{v}\cdot\nabla)\mathbf{v}\bigr) = -\nabla p + \mu\nabla^2\mathbf{v} - \kappa\mu_\mathrm{V1}\ell_P^2\nabla^4\mathbf{v} + (\nabla\times\mathbf{B})\times\mathbf{B} + \rho\mathbf{f}^\mathrm{ext}.\]

The generalized resistive Ohm's law is

\[\mathbf{j} = \sigma\bigl(\mathbf{E} + \mathbf{v}\times\mathbf{B}\bigr).\]

Its \(\mathbf{v}\times\mathbf{B}\) piece is a transport-kinematic
component classified alongside the induction kinematic term.

This eleven-item content set covers every term in standard
incompressible resistive MHD. The classification that follows is over
these eleven items.

\begin{center}\rule{0.5\linewidth}{0.5pt}\end{center}

\subsection{C.4 Architectural Classification of MHD
Terms}\label{c.4-architectural-classification-of-mhd-terms}

The four sub-arc memos NS-MHD-1 through NS-MHD-5 jointly classify every
MHD content channel. Three classes emerge.

\subsubsection{C.4.1 Canonical ED content (6
items)}\label{c.4.1-canonical-ed-content-6-items}

Each item is FORCED at form level by a canonical ED channel and is
ontologically a force sourced by a stable participation structure:

\textbf{(a) Viscous diffusion \(\mu\nabla^2\mathbf{v}\).} The mobility
channel applied component-wise to velocity. Per the field-type-agnostic
vector-extension of the canonical PDE (Architectural Canon Vector
Extension memo), the mobility channel produces a Laplacian-class
diffusion structure on any vector field at fluid scale. Ontologically,
\(\mu\nabla^2\mathbf{v}\) is a participation-flow bias toward smoother
velocity orientation, sourced by the directional-field structure of
\(\mathbf{v}\) itself. Form FORCED; coefficient \(\mu\) INHERITED at
value layer.

\textbf{(b) Magnetic diffusion \(\eta\nabla^2\mathbf{B}\).} The mobility
channel applied component-wise to the magnetic field. Structurally
identical to (a) at the level of canonical-channel derivation --- the
vector-extension is field-type-agnostic, so the same mobility-channel
argument that produces \(\mu\nabla^2\mathbf{v}\) produces
\(\eta\nabla^2\mathbf{B}\). Ontologically, a participation-flow bias on
the magnetic directional field. Form FORCED; resistivity
\(\eta = 1/(\sigma\mu_0)\) INHERITED.

\textbf{(c) Lorentz force \((\nabla\times\mathbf{B})\times\mathbf{B}\).}
Theorem 17 (Gauge-Fields-as-Rule-Type) FORCES the substrate-level
minimal-coupling form \(\partial_\mu \to \partial_\mu - iqA_\mu\) on
charged structural rule-types. The classical semiclassical limit yields
the per-chain force \(q(\mathbf{E} + \mathbf{v}\times\mathbf{B})\);
coarse-graining to fluid scale yields the force density
\(\rho_q\mathbf{E} + \mathbf{j}\times\mathbf{B}\), which reduces in the
MHD limit to \((\nabla\times\mathbf{B})\times\mathbf{B}\).
Ontologically, the bias on charged-chain participation flow imposed by
the gauge directional-field \(\tau_g\). Form FORCED by T17 minimal
coupling; the kinematic \(\mathbf{v}\times\mathbf{B}\) component is
\emph{derived} from the minimal-coupling form rather than separately
committed (NS-MHD-2 §3.2). Coarse-graining bridge INHERITED in parallel
to the NS-2 viscous-content coarse-graining.

\textbf{(d) Maxwell field structure (\(\partial_t\mathbf{B}\),
\(\nabla\times\mathbf{B}\), \(\nabla\cdot\mathbf{B} = 0\)).} Theorem
17's gauge content directly FORCES the four-channel structure of
\(\tau_g\) (group, vertex, worldline, vacuum), from which Maxwell's
equations follow under gauge-quotient identification. The
magnetic-divergence-free constraint \(\nabla\cdot\mathbf{B} = 0\) is the
no-monopoles structural feature of the canonical \(U(1)\) gauge content
--- distinct in status from the velocity incompressibility constraint
(which is fluid-mechanical-imposed). Ontologically, directional-field
dynamics for \(\tau_g\). Form FORCED; field amplitudes INHERITED.

\textbf{(e) R1 substrate-cutoff term
\(-\kappa\mu_\mathrm{V1}\ell_P^2\nabla^4\mathbf{v}\).} The form-FORCED
hyperviscous stabilization arising from V1's finite-width vacuum kernel
(Section §5 of the main paper). Inherits into the MHD momentum equation
by virtue of being a canonical-ED augmentation of the velocity sector
independent of EM coupling. Ontologically, a substrate-level
participation-stability mechanism (T18-class kernel retardation
expressed at V1 vacuum-kernel scale) imposing a hyperviscous bias on
velocity-gradient cascade. Form FORCED; coefficient \(\kappa\)
INHERITED.

All six items are forces in the ED-I-06 sense --- biases on
participation flow sourced by stable participation structures.

\subsubsection{C.4.2 Continuum constraints (non-obstructive, 2
items)}\label{c.4.2-continuum-constraints-non-obstructive-2-items}

These items are not forces in the ED-I-06 sense and are not derivable
from any canonical ED channel. They are continuum-level structural
commitments imposed at fluid scale rather than sourced by participation
structure. Crucially, they are also \emph{non-obstructive}: they
introduce no indefinite-sign Lyapunov contributions and no spectral
pathology.

\textbf{(f) Pressure \(-\nabla p\).} Scalar-gradient in form (the
gradient of a scalar density-class quantity) but \emph{imposed} as a
Lagrange multiplier for the incompressibility constraint, not sourced as
a participation-density bias from below. Ontologically, a
constraint-derived field rather than a force-from-stability.
Sign-definite contribution to the kinetic-energy budget.

\textbf{(g) Velocity incompressibility \(\nabla\cdot\mathbf{v} = 0\).}
Continuum kinematic constraint, not derivable from any canonical
channel. The canonical ED PDE does not commit to divergence-free
participation-measure evolution; the assumption is imposed at the fluid
level. Non-obstructive.

(The magnetic-divergence-free constraint \(\nabla\cdot\mathbf{B} = 0\)
is in a different category, as noted in §C.4.1(d); it is FORCED by
Maxwell field structure under T17 and is canonical ED.)

\subsubsection{C.4.3 Transport-kinematic non-ED terms (obstructive, 3
items)}\label{c.4.3-transport-kinematic-non-ed-terms-obstructive-3-items}

The three remaining items share a single structural class. They are not
sourced by participation structures, are not produced by minimal
coupling, and are dynamically and spectrally obstructive in parallel
ways.

\textbf{(h) Advection \((\mathbf{v}\cdot\nabla)\mathbf{v}\).} Bilinear
in velocity with a directional derivative. Three-angle convergence on
non-ED status established in NS-2.08 (architectural --- no canonical
channel produces this index-and-derivative structure), NS-Smoothness §5
(dynamical --- supplies the unique indefinite-sign contribution to the
gradient-norm Lyapunov via vortex-stretching), and NS-Turbulence §6
(spectral --- bilinear-with-transverse-projector triad incompatible with
P7 symmetric quadratic class).

\textbf{(i) Induction kinematic
\(\nabla\times(\mathbf{v}\times\mathbf{B})\).} Bilinear in
\((\mathbf{v},\mathbf{B})\) with cross-product followed by curl.
Expanded under incompressibility,
\((\mathbf{B}\cdot\nabla)\mathbf{v} - (\mathbf{v}\cdot\nabla)\mathbf{B}\)
--- the magnetic-tension stretching term and the magnetic advection
term, both transport-kinematic bilinear couplings. Three-angle
convergence on non-ED status established in NS-MHD-3: architectural (no
canonical channel for bilinear \(\mathbf{v}\)-\(\mathbf{B}\)
transport-kinematic structure), dynamical (unique indefinite-sign
contribution to magnetic-energy Lyapunov
\(E_M = \tfrac{1}{2}\int|\mathbf{B}|^2\,dV\), structurally parallel to
vortex-stretching's role in 3D NS), and spectral
(bilinear-with-Levi-Civita-projection triad incompatible with P7
symmetric quadratic class).

\textbf{(j) Ohm kinematic \(\mathbf{v}\times\mathbf{B}\).} Cross-product
velocity-magnetic coupling appearing in the generalized Ohm's law. Same
structural class as (h) and (i): bilinear with antisymmetric Levi-Civita
projection between two transport fields, not produced by any canonical
channel.

\textbf{Common structural features.}

\begin{itemize}
\tightlist
\item
  \emph{Not sourced by participation structures.} None corresponds to a
  directional-field bias, a scalar-density gradient, or a
  curvature-like-field bias on participation flow.
\item
  \emph{Frame-kinematic.} All three are bilinear couplings between two
  transport fields that arise specifically when dynamics are written in
  the fluid (Eulerian) coordinate system. They are coordinate artifacts
  of the fluid frame.
\item
  \emph{Same algebraic class.} Bilinear-with-projection between two
  transport fields (transverse projector \(P_{ij}\) for advection;
  Levi-Civita \(\varepsilon_{ijk}\) for induction kinematic and Ohm
  kinematic).
\item
  \emph{Spectrally incompatible with P7.} P7 generates
  symmetric-quadratic triad couplings; the transport-kinematic class is
  antisymmetric / projector-mediated and is structurally incompatible
  with the only canonical nonlinear-coupling channel.
\item
  \emph{Dynamically obstructive.} Each supplies an indefinite-sign
  contribution to a Lyapunov functional that would otherwise be
  monotonically decaying.
\end{itemize}

Ontologically under ED-I-06, these are \emph{frame artifacts} of the
fluid coordinate system --- not forces at all.

\begin{center}\rule{0.5\linewidth}{0.5pt}\end{center}

\subsection{C.5 Final Architectural Classification
Table}\label{c.5-final-architectural-classification-table}

{\def\LTcaptype{none} % do not increment counter
\begin{longtable}[]{@{}
  >{\raggedright\arraybackslash}p{(\linewidth - 10\tabcolsep) * \real{0.1667}}
  >{\raggedright\arraybackslash}p{(\linewidth - 10\tabcolsep) * \real{0.1667}}
  >{\raggedright\arraybackslash}p{(\linewidth - 10\tabcolsep) * \real{0.1667}}
  >{\raggedright\arraybackslash}p{(\linewidth - 10\tabcolsep) * \real{0.1667}}
  >{\raggedright\arraybackslash}p{(\linewidth - 10\tabcolsep) * \real{0.1667}}
  >{\raggedright\arraybackslash}p{(\linewidth - 10\tabcolsep) * \real{0.1667}}@{}}
\toprule\noalign{}
\begin{minipage}[b]{\linewidth}\raggedright
Term
\end{minipage} & \begin{minipage}[b]{\linewidth}\raggedright
Equation
\end{minipage} & \begin{minipage}[b]{\linewidth}\raggedright
Origin
\end{minipage} & \begin{minipage}[b]{\linewidth}\raggedright
Architectural Status
\end{minipage} & \begin{minipage}[b]{\linewidth}\raggedright
Ontological Status (ED-I-06)
\end{minipage} & \begin{minipage}[b]{\linewidth}\raggedright
Notes
\end{minipage} \\
\midrule\noalign{}
\endhead
\bottomrule\noalign{}
\endlastfoot
Viscous diffusion & \(\mu\nabla^2\mathbf{v}\) & Mobility channel on
\(\mathbf{v}\) & Canonical ED & Force from directional field
\(\mathbf{v}\) & Form FORCED; coefficient INHERITED \\
Magnetic diffusion & \(\eta\nabla^2\mathbf{B}\) & Mobility channel on
\(\mathbf{B}\) & Canonical ED & Force from directional field
\(\mathbf{B}\) & Field-type-agnostic vector extension \\
Lorentz force & \((\nabla\times\mathbf{B})\times\mathbf{B}\) & T17
minimal coupling & Canonical ED & Force from directional field
\(\tau_g\) via minimal coupling & Kinematic
\(\mathbf{v}\times\mathbf{B}\) derived from minimal-coupling form \\
R1 substrate cutoff &
\(-\kappa\mu_\mathrm{V1}\ell_P^2\nabla^4\mathbf{v}\) & V1 finite-width
vacuum kernel & Canonical ED & Force from substrate participation
structure & Hyperviscous Lyapunov stabilization \\
Magnetic time-evolution & \(\partial_t\mathbf{B}\) & T17 gauge-field
dynamics & Canonical ED & Directional-field dynamics & Standard
time-derivative on \(A_\mu\) \\
Magnetic divergence-free & \(\nabla\cdot\mathbf{B} = 0\) & T17 / Maxwell
structure & Canonical ED & Directional-field structural constraint &
FORCED by gauge content \\
Pressure & \(-\nabla p\) & Lagrange multiplier for
\(\nabla\cdot\mathbf{v}=0\) & Continuum constraint & Imposed (not a
force) & Non-obstructive; sign-definite \\
Velocity incompressibility & \(\nabla\cdot\mathbf{v} = 0\) & Fluid
kinematic constraint & Continuum constraint & Imposed (not a force) &
Non-obstructive \\
Advection & \((\mathbf{v}\cdot\nabla)\mathbf{v}\) & Eulerian-frame
bilinear & Non-ED (obstructive) & Frame-kinematic (not a force) &
NS-2.08 / NS-Smoothness §5 / NS-Turbulence §6 \\
Induction kinematic & \(\nabla\times(\mathbf{v}\times\mathbf{B})\) &
Eulerian-frame bilinear & Non-ED (obstructive) & Frame-kinematic (not a
force) & NS-MHD-3 three-angle convergence \\
Ohm kinematic & \(\mathbf{v}\times\mathbf{B}\) in Ohm's law &
Eulerian-frame bilinear & Non-ED (obstructive) & Frame-kinematic (not a
force) & Same structural class \\
\end{longtable}
}

\textbf{Aggregate counts.} Of eleven content items: six canonical-ED
forces, two continuum-imposed constraints, three transport-kinematic
non-forces. Only six of eleven items are forces in the ED-I-06 sense;
the remaining five are constraints or frame artifacts of the fluid
coordinate system.

\begin{center}\rule{0.5\linewidth}{0.5pt}\end{center}

\subsection{C.6 Structural Verdict}\label{c.6-structural-verdict}

The four sub-arc closures (NS-MHD-1 / 2 / 3 / 4) and the consolidated
classification table (§C.5) jointly support five structural verdicts:

\textbf{(V1) MHD is partially ED-architectural, exactly like NS.} The
canonical / non-canonical boundary in incompressible MHD is the same
boundary that runs through pure NS: forces from participation structures
on one side, frame-kinematic terms and continuum constraints on the
other. The structural shape of the classification is preserved under the
NS → MHD upgrade.

\textbf{(V2) MHD has strictly more canonical content on the EM side.}
Maxwell field structure, Lorentz force, magnetic divergence-free
constraint, and magnetic diffusion are all canonical ED via Theorem 17
and the mobility channel. The MHD upgrade does not add any new
fluid-mechanical-addition or non-ED obstructive term on the EM side; it
adds only canonical-ED forces. The canonical : non-canonical ratio for
the EM-side additions is 4 : 2; for the pure-NS momentum equation it is
1 : 3. The aggregate MHD ratio (six canonical : five non-canonical) is
dominated by canonical content.

\textbf{(V3) The canonical / non-canonical boundary is the
kinematic-coupling pattern.} Velocity-dependence in continuum equations
splits into two structurally distinct classes:

\begin{itemize}
\tightlist
\item
  \emph{Minimal-coupling-derived velocity-dependence} (Lorentz
  \(\mathbf{v}\times\mathbf{B}\)): canonical ED via Theorem 17. A force
  sourced by the gauge directional-field \(\tau_g\) acting on charged
  chains.
\item
  \emph{Transport-kinematic velocity-dependence} (advection, induction
  kinematic, Ohm kinematic): non-ED. A frame-kinematic artifact of the
  fluid (Eulerian) coordinate system, not a force.
\end{itemize}

This is the load-bearing architectural insight of the MHD arc and
refines the NS-2.08 classification of pure NS by giving it an
ontological grounding under ED-I-06.

\textbf{(V4) The induction term is the magnetic-side analogue of NS
advection.} Both are bilinear-with-projection couplings between two
transport fields. Both are non-ED at architectural / dynamical /
spectral levels. Both supply the unique indefinite-sign contribution to
their respective Lyapunov functionals --- gradient-norm Lyapunov for NS
advection (Section §5 of the main paper), magnetic-energy Lyapunov for
the induction kinematic term (NS-MHD-3 §4.2). The structural parallel is
exact, validating the kinematic-coupling pattern as a generalizable
architectural feature rather than an NS-specific finding.

\textbf{(V5) MHD inherits the same structural obstruction class as NS.}
ED supplies real architectural regularization (R1 + magnetic diffusion +
viscous diffusion), genuine canonical force-content (Lorentz + Maxwell),
and an architectural account of the EM-side additions. But the
obstruction to ED-style smoothness in MHD lies entirely in the
transport-kinematic sector --- advection plus induction kinematic plus
Ohm kinematic --- exactly as in pure NS. The Clay-NS-relevant
decomposition of the main paper extends to MHD with the obstruction
class enlarged from one item (advection) to three items in the same
structural class.

The structural picture is that \textbf{MHD is \emph{more}
ED-architectural than pure NS} in the precise sense that the EM-side
additions are predominantly canonical, while the new
fluid-mechanical-specific structural commitments belong to a single
class with a single architectural origin (transport-kinematic) and a
single architectural defect (no canonical channel, frame-kinematic in
ED-I-06 ontology, spectrally incompatible with P7).

\begin{center}\rule{0.5\linewidth}{0.5pt}\end{center}

\subsection{C.7 Relation to the Main NS
Synthesis}\label{c.7-relation-to-the-main-ns-synthesis}

This appendix extends the architectural picture of the main paper
without revising any of its conclusions:

\begin{itemize}
\tightlist
\item
  The main paper's Section §3 (NS-1 dimensional forcing via Path
  B-strong) is preserved.
\item
  The main paper's Section §4 (NS-2 form derivation via partial
  vector-extension) extends to MHD via the field-type-agnostic
  vector-extension argument (§C.4.1(b), magnetic diffusion).
\item
  The main paper's Section §5 (NS-Smoothness Intermediate Path C
  verdict, R1 mechanism, vortex-stretching as unique indefinite-sign
  Lyapunov contribution) extends to MHD with the induction kinematic
  term playing the same structural role on the magnetic side (§C.6 V4).
\item
  The main paper's Section §6 (NS-Turbulence H1 trivial / H2 partial /
  H3 fail; advection spectrally incompatible with P7) extends to MHD
  with induction kinematic and Ohm kinematic spectrally incompatible
  with P7 in the same way.
\item
  The main paper's three-angle convergence on advection-as-non-ED
  (architectural NS-2.08, dynamical NS-Smoothness, spectral
  NS-Turbulence) is now joined by a second three-angle convergence on
  the induction kinematic term (NS-MHD-3), validating the
  kinematic-coupling pattern as a generalizable architectural feature.
\end{itemize}

The ED-I-06 ontology supplies the conceptual unification under which all
of these structural results read as instances of a single
forces-vs-frame-kinematic-or-constraint distinction. The canonical /
non-canonical boundary, established in the main paper at
structural-architectural level, is now also grounded ontologically:
canonical ED content is force-from-stability; non-ED content is
frame-kinematic or imposed-constraint.

The main paper's headline disclaimers remain unchanged: ED does not
resolve whether 3D NS blows up at finite time, does not predict critical
Reynolds numbers, and does not supply a turbulence-cascade architectural
template. The MHD extension preserves these disclaimers and adds
parallel ones for MHD: ED does not derive dynamo theory, does not supply
an MHD-turbulence cascade architectural template, and does not predict
the magnetic-reconnection rate. What ED \emph{does} deliver --- for both
NS and MHD, under both structural and ontological readings --- is a
complete account of which content channels are forces sourced by
participation structures, which are frame artifacts of the fluid
coordinate system, and which are continuum-imposed constraints. That
account is unchanged in substance by the integration of the ED-I-06
ontology; it is strengthened in framing.

\begin{center}\rule{0.5\linewidth}{0.5pt}\end{center}

\emph{Appendix C closes the integration of the NS-MHD arc and the
ED-I-06 ontology into the NS Synthesis Paper. Eleven incompressible-MHD
content items are classified into six canonical-ED forces, two
continuum-imposed constraints, and three transport-kinematic non-forces.
The kinematic-coupling pattern --- minimal-coupling-derived
velocity-dependence canonical, transport-kinematic velocity-dependence
non-ED --- is the canonical / non-canonical boundary in continuum fluid
mechanics. The ED-I-06 ontology grounds this boundary as the
forces-vs-frame-kinematic distinction, unifying the classification
results of NS-2.08, NS-Smoothness, NS-Turbulence, NS-MHD-2, NS-MHD-3,
and NS-MHD-4 under a single ontological framework. The main NS Synthesis
Paper's results are preserved and ontologically strengthened.}

\begin{center}\rule{0.5\linewidth}{0.5pt}\end{center}

\section{Appendix D --- Diffusion Coarse-Graining Theorem
(DCGT)}\label{appendix-d-diffusion-coarse-graining-theorem-dcgt}

\emph{Companion appendix to} The Architectural Foundations of
Navier-Stokes in Event Density: Form Derivation, Clay-Relevance
Decomposition, and the Structural Status of Turbulence.

\begin{center}\rule{0.5\linewidth}{0.5pt}\end{center}

\subsection{D.1 Purpose}\label{d.1-purpose}

This appendix presents the \textbf{Diffusion Coarse-Graining Theorem
(DCGT)} proved across the six memos of Arc D in the NS / MHD program.
The theorem closes the last INHERITED bridges in the architectural
classification of the main paper and Appendix C, and completes the
substrate-to-continuum foundation of the NS Synthesis Paper.

Specifically, this appendix:

\begin{itemize}
\tightlist
\item
  \textbf{Presents the Diffusion Coarse-Graining Theorem} in canonical
  publication form, parallel in structure to Theorem 17
  (Gauge-Fields-as-Rule-Type), Theorem 18 (V1 Kernel Retardation), and
  Theorem 19 (Newton's Law from Substrate).
\item
  \textbf{Shows how scalar diffusion, viscosity, R1 hyperviscosity,
  viscoelastic memory, and the Lorentz force all arise from substrate
  primitives} under hydrodynamic-window coarse-graining as outputs of a
  single multi-scale expansion. Five canonical-ED dynamical content
  channels of NS / MHD share a single substrate origin.
\item
  \textbf{Closes the last INHERITED bridges in NS-2 and NS-MHD-2.} The
  form FORCED / coarse-graining INHERITED status of NS-2 viscous content
  and NS-MHD-2 Lorentz force is promoted to form FORCED /
  coarse-graining FORCED-via-DCGT.
\item
  \textbf{Completes the substrate-to-continuum foundation of the NS
  Synthesis Paper.} With DCGT in place, every canonical-ED dynamical
  term in the NS / MHD architectural classification is
  substrate-grounded; the paper transitions from
  architecturally-complete-with-INHERITED-coarse-graining to
  architecturally-and-substrate-complete.
\end{itemize}

The theorem is FORCED-unconditional. It joins Theorems 17 / 18 / 19 in
the structural-foundation inventory of the Event Density program.

\begin{center}\rule{0.5\linewidth}{0.5pt}\end{center}

\subsection{D.2 Hydrodynamic Window and Coarse-Graining
Operator}\label{d.2-hydrodynamic-window-and-coarse-graining-operator}

\textbf{Hydrodynamic window.} A coarse-graining cell
\(\Omega(\mathbf{x},R_\mathrm{cg})\) of radius \(R_\mathrm{cg}\)
centered at macroscopic position \(\mathbf{x}\) is in the
\emph{hydrodynamic window} iff three scale-separation conditions hold:

\[\ell_P \;\ll\; R_\mathrm{cg} \;\ll\; L_\mathrm{flow},\]

where \(\ell_P\) is the substrate length scale (V1 kernel width) and
\(L_\mathrm{flow}\) is the macroscopic flow gradient scale. The lower
bound suppresses substrate-discreteness fluctuations
(statistical-regularity condition); the upper bound preserves field
gradients (gradient-expansion-truncation validity). For temporal kernels
(V5), the analogous temporal window is
\(\tau_\mathrm{V5} \ll \tau_\mathrm{cg} \ll \tau_\mathrm{flow}\).

\textbf{Coarse-graining operator.} Define the spatial coarse-graining
operator on a substrate field \(f\) as

\[\langle f\rangle_{R_\mathrm{cg}}(\mathbf{x}) \;\equiv\; \frac{1}{|\Omega(\mathbf{x},R_\mathrm{cg})|}\int_{\Omega(\mathbf{x},R_\mathrm{cg})} f(\mathbf{x}+\boldsymbol{\delta})\,d^3\boldsymbol{\delta}.\]

The temporal analogue, used for V5-mediated cross-chain memory
contributions, is

\[\langle f\rangle^\mathrm{(t)}_{\tau_\mathrm{cg}}(\mathbf{x},t) \;\equiv\; \int_0^\infty K_\mathrm{cg}(\tau)\,f(\mathbf{x},t-\tau)\,d\tau.\]

\textbf{Multi-scale expansion principle.} All continuum operators in the
DCGT arise from the multi-scale expansion of the coarse-graining
operator applied to substrate fluxes --- event flux \(\mathbf{J}_\rho\)
for the scalar sector, momentum-flux tensor \(\Pi_{ij}\) for the
directional-field sector, charged-chain momentum-flux \(\Pi^{(q)}_{ij}\)
for the EM coupling, V5-mediated temporal kernel for viscoelastic
memory. Even moments of V1 produce successive Laplacian-class
corrections (zeroth: field itself; second: leading diffusion /
viscosity; fourth: R1 hyperviscosity); odd moments vanish by V1
isotropy; V5 temporal moments produce Maxwell-class memory
contributions; T17 minimal coupling shifts the charged-chain momentum
and produces the Lorentz force at the divergence.

A single coarse-graining operator generates every canonical-ED continuum
dynamical content channel.

\begin{center}\rule{0.5\linewidth}{0.5pt}\end{center}

\subsection{D.3 Statement of the Diffusion Coarse-Graining Theorem
(DCGT)}\label{d.3-statement-of-the-diffusion-coarse-graining-theorem-dcgt}

\begin{center}\rule{0.5\linewidth}{0.5pt}\end{center}

\textbf{Theorem (Diffusion Coarse-Graining Theorem,
FORCED-unconditional).}

\begin{quote}
\emph{Let the ED substrate consist of chains carrying event density
\(\rho\), velocity \(\mathbf{v}\), and (for charged-chain populations)
charge density \(\rho_q\), evolving under V1 (finite-width spatial
vacuum kernel, Theorem 18) and V5 (finite-memory cross-chain temporal
kernel) participation kernels with substrate-level minimal coupling per
Theorem 17 in the gauge sector.}

\emph{Under coarse-graining \(\langle\,\cdot\,\rangle_{R_\mathrm{cg}}\)
over a hydrodynamic window}

\[\ell_P \;\ll\; R_\mathrm{cg} \;\ll\; L_\mathrm{flow},\]

\emph{the continuum-scale evolution equations for scalar and directional
fields are:}
\end{quote}

\[\boxed{\;\begin{aligned}
\partial_t\rho \;&=\; \nabla\cdot\bigl(M(\rho)\,\nabla\rho\bigr), \\[4pt]
\rho\,\partial_t\mathbf{v} \;&=\; \mu(\rho)\,\nabla^2\mathbf{v} \;-\; \kappa\,\mu_\mathrm{V1}\,\ell_P^2\,\nabla^4\mathbf{v} \;+\; \lambda(\rho)\,\partial_t\mathbf{v} \\[2pt]
&\qquad +\; \rho_q\,\mathbf{E} \;+\; \mathbf{j}\times\mathbf{B} \;-\; \rho(\mathbf{v}\cdot\nabla)\mathbf{v} \;-\; \nabla p,
\end{aligned}\;}\]

\begin{quote}
\emph{with truncated higher-order multi-scale-expansion remainder
bounded by}

\[\mathcal{O}\!\left(\bigl(\ell_P/L_\mathrm{flow}\bigr)^4\right) \;+\; \mathcal{O}\!\left(\bigl(\tau_\mathrm{V5}/\tau_\mathrm{flow}\bigr)^2\right).\]
\end{quote}

\textbf{Properties.}

\textbf{(1) Form-FORCED.} All continuum operators ---
\(\nabla\cdot(M\nabla)\), \(\mu\nabla^2\), \(-\kappa\ell_P^2\nabla^4\),
\(\lambda\partial_t\), \(\rho_q\mathbf{E} + \mathbf{j}\times\mathbf{B}\)
--- arise unconditionally from the multi-scale expansion of substrate
flux statistics under hydrodynamic-window coarse-graining. No
fluid-mechanical postulate enters the derivation.

\textbf{(2) Sign-FORCED.} The V1 kernel's positive Fourier transform
(Theorem N1 admissibility class, Theorem 18 forward-cone support) forces
positive even moments. Hence the diffusion / viscosity coefficients
\(M\), \(\mu\) are non-negative; the R1 coefficient
\(\kappa\mu_\mathrm{V1}\ell_P^2\) is positive (R1 enters the equation of
motion with a stabilizing negative sign); the V5 memory coefficient
\(\lambda\) is non-negative.

\textbf{(3) Value-INHERITED.} All numerical coefficients ---
\(M(\rho)\), \(\mu(\rho)\), \(\mu_\mathrm{V1}\), \(\kappa\),
\(\lambda(\rho)\), \(\rho_\mathrm{max}\),
\(\bigl\langle\delta^2\bigr\rangle_{V1}\),
\(\bigl\langle\delta^4\bigr\rangle_{V1}\),
\(\bigl\langle\tau\bigr\rangle_{V5}\) --- are INHERITED at value layer
from the V1 / V5 kernel profiles and the
participation-bandwidth-modulated transition rate \(\Gamma_0(\rho)\).
DCGT fixes their structural form (form FORCED) but not their pointwise
numerical values, in keeping with the program's standard form-FORCED /
value-INHERITED methodology.

\textbf{(4) Kinematic-Coupling Pattern (Substrate-Confirmed).}
Transport-kinematic terms --- advection
\((\mathbf{v}\cdot\nabla)\mathbf{v}\) in NS, induction kinematic
\(\nabla\times(\mathbf{v}\times\mathbf{B})\) in MHD, Ohm-law kinematic
component \(\mathbf{v}\times\mathbf{B}\) --- do \emph{not} arise from
substrate force-from-stability mechanisms. They appear in the continuum
equations as Eulerian-frame bookkeeping of convective fluxes
(\(\partial_j(\rho v_i v_j)\) for advection; analogous projections for
induction and Ohm). They preserve their non-ED transport-kinematic
frame-artifact status from the main paper Sections §3--§6 and Appendix
C, now substrate-confirmed via the substrate derivation. The
kinematic-coupling pattern (minimal-coupling-derived velocity-dependence
canonical / transport-kinematic velocity-dependence non-ED) is
substrate-grounded.

\begin{center}\rule{0.5\linewidth}{0.5pt}\end{center}

The boxed system gives the \textbf{substrate-grounded continuum
dynamical equations of NS / MHD}. Every canonical-ED term has substrate
origin via DCGT; every non-ED term preserves its frame-kinematic or
constraint status from Appendix C.

\begin{center}\rule{0.5\linewidth}{0.5pt}\end{center}

\subsection{D.4 Proof Sketch}\label{d.4-proof-sketch}

The proof aggregates five parallel applications of the same
multi-scale-expansion argument.

\subsubsection{(1) Scalar diffusion}\label{scalar-diffusion}

Cell-averaged participation density
\(\rho(\mathbf{x},t) = N_\Omega/|\Omega|\) satisfies local conservation
\(\partial_t\rho = -\nabla\cdot\mathbf{J}_\rho\). V1-mediated chain-step
transition statistics with isotropic kernel + first-order gradient
expansion of \(\Gamma_0(\rho)\rho\) across the cell boundary:
leading-order term cancels by isotropy of \(K_{V1}\) (forward and
backward contributions match), and the gradient term yields
\(\mathbf{J}_\rho = -M(\rho)\nabla\rho + \mathcal{O}(\ell_P^2\nabla^3\rho)\).
Substituting into the continuity equation produces
\(\partial_t\rho = \nabla\cdot(M(\rho)\nabla\rho) + \mathcal{O}(\ell_P^2\nabla^4\rho)\).
The mobility coefficient
\(M(\rho) \propto \partial[\Gamma_0(\rho)\rho]/\partial\rho \cdot \bigl\langle\delta^2\bigr\rangle_{V1}\)
is identified at substrate level.

\subsubsection{(2) Directional-field diffusion
(viscosity)}\label{directional-field-diffusion-viscosity}

Cell-averaged velocity \(\mathbf{v} = P_\Omega/M_\Omega\) satisfies
local momentum conservation
\(\partial_t(\rho v_i) = -\partial_j\Pi_{ij}\). Chain-step
momentum-transfer statistics with V1 isotropy + first-order gradient
expansion of \(v_i\): V1 second-moment integral diagonal
\(\propto\delta^{jk}\) by isotropy; antisymmetric part of the gradient
cancels under the diagonal contraction; symmetric Newtonian deviatoric
stress \(-\mu(\rho)(\partial_i v_j + \partial_j v_i)\) survives.
Substituting into momentum conservation and using continuity yields
\(\rho\,\partial_t\mathbf{v} = \mu(\rho)\nabla^2\mathbf{v} - \rho(\mathbf{v}\cdot\nabla)\mathbf{v} - \nabla p + \mathcal{O}(\ell_P^2\nabla^4\mathbf{v})\).
Substrate-derived viscosity
\(\mu(\rho) = \tfrac{1}{3}\rho\bigl\langle\delta^2\bigr\rangle_{V1}\Gamma_0(\rho)\).

\subsubsection{(3) R1 hyperviscosity}\label{r1-hyperviscosity}

Continue the V1-weighted multi-scale expansion to the next non-vanishing
moment. The V1 fourth-moment \(\bigl\langle\delta^4\bigr\rangle_{V1}\)
contributes a \(\nabla^3\mathbf{v}\) term to the momentum flux.
Momentum-conservation divergence of this term yields
\(-\kappa\,\mu_\mathrm{V1}\,\ell_P^2\nabla^4\mathbf{v}\) in the equation
of motion, with \(\kappa\) a kurtosis-class profile coefficient (ratio
of fourth to squared second moment) and \(\mu_\mathrm{V1}\) the V1
amplitude prefactor. Sign positive by positive-Fourier-transform
property of V1; the resulting term enters the equation of motion with
stabilizing negative sign. The substrate derivation matches the
form-FORCED top-down derivation of Section §5 of the main paper exactly:
same operator form (\(\nabla^4\mathbf{v}\)), same sign (stabilizing),
same \(\ell_P^2\) scaling, same physical origin (V1 finite-width vacuum
kernel).

\subsubsection{(4) Viscoelastic memory}\label{viscoelastic-memory}

V5 cross-chain participation kernel acts on velocity history via
temporal convolution. Hydrodynamic-window temporal coarse-graining
\(\bigl[K_{V5}\ast\partial_t\mathbf{v}\bigr] = \partial_t\mathbf{v} - \bigl\langle\tau\bigr\rangle_{V5}\partial_t^2\mathbf{v} + \cdots\)
generates the Maxwell-class memory ODE
\(\tau_R\dot\sigma + \sigma = 2\mu S\) with
\(\tau_R = \bigl\langle\tau\bigr\rangle_{V5}\) identified as the V5
first temporal moment. At hydrodynamic flow scales, the leading effect
renormalizes the inertia: a \(\lambda(\rho)\partial_t\mathbf{v}\)
contribution to the momentum equation with
\(\lambda(\rho) \propto \rho\,\tau_\mathrm{V5}\,\Gamma_0(\rho)\).

\subsubsection{(5) Lorentz force}\label{lorentz-force}

Theorem 17 minimal coupling \(\partial_\mu \to \partial_\mu - iqA_\mu\)
on charged chains shifts canonical momentum \(p_i \to p_i + qA_i\).
Charged-chain momentum-flux tensor picks up minimal-coupling
contribution \(\delta\Pi^{(q)}_{ij} = j_j A_i\) (with
\(\mathbf{j} = \rho_q\mathbf{v}\)). Momentum-conservation divergence:
\(-\partial_j(j_jA_i) = -j_j\partial_jA_i\) (the \(-A_i\partial_j j_j\)
piece eliminated by charge conservation). Cross-product identity
\(\mathbf{j}\times\mathbf{B} = \mathbf{j}\times(\nabla\times\mathbf{A})\)
decomposes
\(-j_j\partial_jA_i = [\mathbf{j}\times\mathbf{B}]_i - j_j\partial_iA_j\).
Time-component minimal-coupling contribution \(-\partial_t(\rho_qA_i)\)
combines with scalar-potential gradient \(-\rho_q\partial_i\phi\) via
\(\mathbf{E} = -\nabla\phi - \partial_t\mathbf{A}\) to yield
\(\rho_qE_i\). Combining:
\(\rho\partial_t\mathbf{v} \supset \rho_q\mathbf{E} + \mathbf{j}\times\mathbf{B}\).
Force-from-participation-structure status confirmed; no
transport-kinematic structure appears.

\subsubsection{Aggregation}\label{aggregation}

Each of (1)--(5) is a separate but parallel application of the same
coarse-graining operator + multi-scale expansion. The continuum
dynamical equations of the boxed system are obtained by inserting all
substrate-derived flux contributions into the local conservation laws
and identifying the resulting continuum operators by their
multi-scale-expansion order. The error bounds
\(\mathcal{O}((\ell_P/L_\mathrm{flow})^4) + \mathcal{O}((\tau_\mathrm{V5}/\tau_\mathrm{flow})^2)\)
follow directly from the truncated higher-order V1 / V5 kernel moments.

The non-ED transport-kinematic and continuum-constraint terms appear in
the boxed system not as outputs of the coarse-graining derivation but as
Eulerian-bookkeeping artifacts (advection) and continuum-level
structural commitments (pressure, incompressibility) preserved from
Appendix C. DCGT does not derive them; it confirms their non-ED status
by \emph{exhibiting} the canonical-ED side from substrate primitives and
showing that the transport-kinematic / constraint side does not emerge.

The structural unity of the proof is the coarse-graining operator's
universality: a single operator, applied to substrate fluxes and
substrate kernels, generates every canonical-ED continuum dynamical
content channel of NS / MHD.

\begin{center}\rule{0.5\linewidth}{0.5pt}\end{center}

\subsection{D.5 Structural Consequences for the NS
Program}\label{d.5-structural-consequences-for-the-ns-program}

DCGT closure has six direct structural consequences for the NS / MHD
program:

\textbf{(i) NS-2 viscous content is fully substrate-derived.} Section §4
of the main paper derives the standard Newtonian-fluid NS form via
partial vector-extension of the canonical PDE plus chain-substrate
coarse-graining heuristics. The coarse-graining step was form FORCED /
INHERITED at the level of the bridge derivation. DCGT closes the bridge:
chain momentum-flux statistics under V1-mediated transitions,
coarse-grained over the hydrodynamic window, produce
\(\mu(\rho)\nabla^2\mathbf{v}\) at fluid scale. The NS-2 status promotes
from form FORCED / coarse-graining INHERITED to form FORCED /
coarse-graining FORCED-via-DCGT.

\textbf{(ii) NS-MHD-2 Lorentz force is fully substrate-derived.}
Appendix C §C.4.1(c) classifies the Lorentz force as canonical ED via
Theorem 17 minimal coupling, with the coarse-graining bridge INHERITED.
DCGT closes the bridge: T17 substrate-level minimal coupling on charged
chains, coarse-grained via the chain-momentum-flux machinery, produces
\(\rho_q\mathbf{E} + \mathbf{j}\times\mathbf{B}\) at fluid scale. The H2
verdict (Lorentz force canonical ED) is now substrate-grounded.

\textbf{(iii) R1 is substrate-derived.} The form-FORCED hyperviscous
stabilization \(-\kappa\mu_\mathrm{V1}\ell_P^2\nabla^4\mathbf{v}\)
established in Section §5 of the main paper is no longer an
architecturally-postulated term arising from V1's finite-width
participation kernel via top-down argument. It is a substrate-derived
next-to-leading-order multi-scale-expansion correction from V1's
fourth-moment / second-moment ratio. The two derivations are reconciled.
Consequently the gradient-norm Lyapunov decay structure
\(-\kappa\mu_\mathrm{V1}\ell_P^2\|\nabla^3\mathbf{v}\|_2^2\) in ED-only
NS is preserved unchanged with substrate grounding.

\textbf{(iv) P4-NN viscoelasticity is substrate-derived.} The
Maxwell-class memory ansatz \(\tau_R\dot\sigma + \sigma = 2\mu S\) used
in the companion P4-NN rheology paper is no longer a postulated
continuum constitutive form; it is derived from V5 cross-chain
temporal-memory at hydrodynamic-window scale, with
\(\tau_R = \bigl\langle\tau\bigr\rangle_{V5}\) identified as the V5
first temporal moment INHERITED at value layer. The companion paper's
viscoelastic content is structurally strengthened.

\textbf{(v) The canonical / non-canonical boundary is
substrate-confirmed.} The classification of NS / MHD content into
canonical-ED forces vs.~transport-kinematic frame artifacts
vs.~continuum constraints (NS-2.08 in the main paper, §C.4 in Appendix
C) is no longer an architectural classification only; it is a
substrate-grounded result. Every canonical-ED term emerges from
substrate primitives via DCGT; transport-kinematic terms emerge as
Eulerian-bookkeeping artifacts at the level of momentum-flux divergence
(the advection term explicitly visible as \(\partial_j(\rho v_iv_j)\) at
the substrate-coarse-graining stage); continuum constraints remain
imposed at fluid scale. The kinematic-coupling pattern
(minimal-coupling-derived velocity-dependence canonical /
transport-kinematic velocity-dependence non-ED) is substrate-confirmed.

\textbf{(vi) The NS Synthesis Paper is now architecturally and
substrate-complete.} Sections §3--§6 of the main paper plus Appendix C
plus the present Appendix D give a complete account at three levels: (a)
architectural --- what canonical-ED operator structure each NS / MHD
term has; (b) ontological (ED-I-06) --- what kind of participation
structure each term sources, or fails to source; (c) substrate --- how
each canonical-ED term emerges from substrate primitives via
hydrodynamic-window coarse-graining. The paper transitions from a
structural classification with INHERITED bridges to a substrate-grounded
classification with all bridges closed.

The structural picture: \textbf{DCGT is the substrate-to-continuum
bridge theorem of the NS / MHD program}, occupying the same load-bearing
role for fluid-mechanical foundations that Theorem 18 occupies for the
kernel-level arrow of time, Theorem 19 for substrate-level Newton's
gravity, and Theorem 17 for gauge-fields-as-rule-type.

\begin{center}\rule{0.5\linewidth}{0.5pt}\end{center}

\subsection{D.6 Relation to the Main
Text}\label{d.6-relation-to-the-main-text}

This appendix interacts with the main paper and Appendix C in five
specific ways:

\textbf{(a) Replaces the last INHERITED caveats in NS-2 and NS-MHD-2.}
The main paper's Section §4 (NS-2 form derivation) and Appendix C §C.4.1
(Lorentz force canonical ED) state form FORCED with INHERITED
coarse-graining bridges. With DCGT, both bridges are closed; the
INHERITED caveats are removed. Future reprints of the synthesis paper
may cite Appendix D in place of the INHERITED-coarse-graining
qualifiers.

\textbf{(b) Strengthens the architectural classification with substrate
grounding.} Appendix C presents the eleven-item architectural
classification of incompressible MHD with canonical / non-canonical
boundary established via three-angle convergence. Appendix D supplies
the substrate-derivation of every canonical-ED entry in that
classification. The two appendices are complementary: Appendix C
establishes the architectural-and-ontological classification; Appendix D
supplies the substrate-derivation foundation.

\textbf{(c) Completes the ED-canonical picture of continuum fluid
mechanics.} Combined with Appendix C's ED-I-06 ontological reading,
Appendix D's substrate derivation, and the main paper's architectural
and Clay-relevance results, the NS Synthesis Paper now presents a
complete ED-canonical account of continuum fluid mechanics: substrate
primitives → DCGT → canonical PDE → NS form → architectural
classification → ED-I-06 ontological reading → Clay-relevance
decomposition → MHD extension. Every layer is now closed.

\textbf{(d) Positions the NS Synthesis Paper as a fully closed arc.} The
pre-DCGT version of the paper had a structural hole at the
coarse-graining bridge --- an honest INHERITED qualifier that any
kinetic-theory-versed reader would notice. With Appendix D, the hole is
closed. The paper is fully closed both architecturally and
substrate-foundationally; no INHERITED bridges remain in the
canonical-ED sector.

\textbf{(e) Preserves the main paper's headline disclaimers.} DCGT does
not resolve whether 3D NS blows up at finite time, does not predict
critical Reynolds numbers, does not supply a turbulence-cascade
architectural template, does not derive dynamo theory or the
magnetic-reconnection rate. These honest disclaimers from the main paper
and Appendix C remain unchanged. What DCGT delivers is the
substrate-grounding of \emph{what ED does account for}: the canonical-ED
dynamical content of NS and MHD. The non-ED content --- advection,
induction kinematic, Ohm kinematic, pressure, incompressibility ---
preserves its non-canonical status, and the obstructions to ED-style
smoothness in 3D NS / MHD remain in the transport-kinematic sector
exactly as the main paper describes.

The NS Synthesis Paper, as of the addition of Appendix D, is
architecturally complete (main paper Sections §3--§6), ontologically
grounded (Appendix C under ED-I-06), MHD-extended (Appendix C
eleven-item classification), and substrate-grounded (Appendix D via
DCGT). The paper closes the NS / MHD architectural-and-substrate program
of the Event Density framework.

\begin{center}\rule{0.5\linewidth}{0.5pt}\end{center}

\emph{Appendix D presents the Diffusion Coarse-Graining Theorem as the
substrate-to-continuum bridge for NS / MHD canonical-ED dynamical
content. Five canonical-ED content channels --- scalar diffusion,
viscosity, R1 hyperviscosity, viscoelastic memory, Lorentz force ---
derived from a single multi-scale expansion of the coarse-graining
operator applied to substrate flux statistics under hydrodynamic-window
scale separation \(\ell_P \ll R_\mathrm{cg} \ll L_\mathrm{flow}\). Form
FORCED, signs FORCED stabilizing by V1 positive-Fourier-transform,
values INHERITED from V1 / V5 kernel profiles. Kinematic-coupling
pattern substrate-confirmed: transport-kinematic terms emerge as
Eulerian-bookkeeping frame-artifacts; continuum constraints preserved as
fluid-mechanical commitments. NS-2 viscous-content + NS-MHD-2 Lorentz
coarse-graining INHERITED bridges closed; R1 substrate origin grounded;
P4-NN viscoelastic ansatz substrate-grounded; canonical / non-canonical
boundary substrate-confirmed. The NS Synthesis Paper is architecturally
and substrate-complete; DCGT joins Theorems 17 / 18 / 19 in the
structural-foundation inventory of the Event Density program.}

\begin{center}\rule{0.5\linewidth}{0.5pt}\end{center}

\section{Appendix E --- Yang-Mills Synthesis and Clay-Relevance
Statement}\label{appendix-e-yang-mills-synthesis-and-clay-relevance-statement}

\emph{Companion appendix to} The Architectural Foundations of
Navier-Stokes in Event Density: Form Derivation, Clay-Relevance
Decomposition, and the Structural Status of Turbulence.

\begin{center}\rule{0.5\linewidth}{0.5pt}\end{center}

\subsection{E.1 Purpose}\label{e.1-purpose}

This appendix extends the architectural and substrate-foundation work of
the main paper and Appendices C and D to the non-Abelian gauge sector.
It summarizes the six-memo Yang-Mills (YM) arc closed in the program on
2026-04-30 and presents the structural results that this arc produced.

Specifically, this appendix:

\begin{itemize}
\tightlist
\item
  \textbf{Summarizes the six-memo Yang-Mills arc} (YM-1 through YM-6)
  and its load-bearing structural findings.
\item
  \textbf{Presents the substrate → continuum Yang-Mills derivation} of
  the canonical YM equation \(D_\mu F^{\mu\nu} = J^\nu\) from ED
  substrate primitives via the same DCGT machinery developed in Appendix
  D for the fluid-mechanical sector.
\item
  \textbf{States the substrate mass-gap mechanism and its conditional
  survival criterion} under the continuum limit, parallel in framing to
  NS-Smoothness's Intermediate Path C verdict on Clay-NS.
\item
  \textbf{Summarizes the architectural classification and OS-positivity
  audit} of YM content under the ED-I-06 ontological reading established
  in Appendix C.
\item
  \textbf{Provides the Clay-relevance statement} for the Yang-Mills
  Existence and Mass Gap Clay Millennium Problem, parallel in form and
  honesty to the main paper's Clay-NS-relevance decomposition.
\end{itemize}

The appendix is publication-grade and stylistically consistent with
Appendices C and D. The main paper's results remain unchanged; the
appendix extends the architectural picture from the fluid-mechanical
sector to the non-Abelian gauge sector, supplying the YM analogue of the
NS / MHD content already classified.

\begin{center}\rule{0.5\linewidth}{0.5pt}\end{center}

\subsection{E.2 Substrate → Continuum Yang-Mills Equation
(YM-2)}\label{e.2-substrate-continuum-yang-mills-equation-ym-2}

The substrate-derivation methodology of Appendix D (DCGT) generalizes
from the Abelian gauge sector (Lorentz force in NS-MHD-2 / Appendix C
§C.4.1(c)) to the non-Abelian sector via T17 generalized minimal
coupling on charged structural rule-types.

\textbf{Substrate-derived continuum YM equation.} Under
hydrodynamic-window coarse-graining
\(\ell_P \ll R_\mathrm{cg} \ll L_\mathrm{flow}\) over ED substrate
primitives (chains, V1 / V5 kernels, charged structural rule-types under
T17 generalized minimal coupling on a compact simple gauge group \(G\)
with Lie algebra \(\mathfrak{g}\)), the cell-averaged non-Abelian
gauge-field strength \(F_{\mu\nu}^a(\mathbf{x},t)\) and gauge-current
density \(J^{\nu\,a}(\mathbf{x},t)\) satisfy:

\[\boxed{\;D_\mu F^{\mu\nu} \;=\; J^\nu, \qquad F_{\mu\nu}^a \;=\; \partial_\mu A_\nu^a - \partial_\nu A_\mu^a + gf^{abc}A_\mu^b A_\nu^c,\;}\]

with covariant derivative \(D_\mu = \partial_\mu - igA_\mu^a T^a\)
acting on the appropriate matter representation, and gauge-current
covariantly conserved \(D_\mu J^\mu = 0\) via T17 gauge-quotient
identification.

\textbf{Three structural identifications.}

\begin{itemize}
\tightlist
\item
  \textbf{Field strength \(F_{\mu\nu}^a\)} arises from substrate-level
  directional-field participation curvature. The non-Abelian commutator
  term \(-ig[A_\mu, A_\nu]\) is FORCED by T17 clause C2 --- the
  substrate gauge group's non-Abelian-capable Lie-bracket structure on
  the generators \([T^a, T^b] = if^{abc}T^c\). Substrate-level
  participation-channel parallel-transport on a closed loop accrues a
  phase \(\sim g[A_\mu, A_\nu]\,\delta^\mu\delta^\nu\) at second order,
  generated by the rule-type bracket; closing and antisymmetrizing
  extracts the substrate field strength.
\item
  \textbf{Covariant derivative \(D_\mu\)} arises from T17 generalized
  minimal coupling at the vertex level:
  \(\partial_\mu \to \partial_\mu - igA_\mu^a T^a\) acting on charged
  structural rule-types. This is the non-Abelian generalization of the
  \(U(1)\) minimal coupling \(\partial_\mu \to \partial_\mu - iqA_\mu\)
  used in Appendix D §D.4 step (5) for the Lorentz force.
\item
  \textbf{Matter source \(J^{\nu\,a}\)} arises from charged-chain flux
  at fluid scale. The non-Abelian generalization of the
  \(\mathbf{j} = \rho_q\mathbf{v}\) current in Appendix D, with the
  gauge-component index \(a\) added to the spacetime current. Charge
  conservation \(D_\mu J^\mu = 0\) is FORCED by T17 clauses C3 + C8
  (vertex-quotient + unified four-channel quotient).
\end{itemize}

\textbf{Status.} Form FORCED by T17 generalized minimal coupling + DCGT
multi-scale expansion + non-Abelian rule-type bracket structure.
Sign-FORCED stabilizing by V1 positive Fourier transform.
Value-INHERITED at coupling \(g\), kernel widths, and specific compact
simple gauge group choice. Bianchi identity \(D_{[\mu}F_{\nu\rho]} = 0\)
preserved as algebraic structural consequence. The Abelian limit
(\(f^{abc} = 0\)) recovers the inhomogeneous Maxwell equation
\(\partial_\mu F^{\mu\nu} = J^\nu\) of Appendix C.

The substrate derivation of the YM equation is structurally parallel to
the substrate derivation of the Lorentz force in Appendix D §D.4 step
(5); the same coarse-graining operator + multi-scale expansion +
minimal-coupling momentum-flux divergence machinery applies to the
non-Abelian case with the addition of the rule-type bracket structure.

\begin{center}\rule{0.5\linewidth}{0.5pt}\end{center}

\subsection{E.3 Substrate Mass-Gap Mechanism
(YM-3)}\label{e.3-substrate-mass-gap-mechanism-ym-3}

The substrate-derived YM equation acquires a substrate-cutoff mass-gap
mechanism via the V1 kernel's finite spatial width.

\textbf{Mechanism.} Four structural elements:

\begin{itemize}
\tightlist
\item
  \textbf{V1 finite width \(\sim \ell_P\).} FORCED at substrate level by
  Theorem N1 (V1 kernel admissibility class) and Theorem 18
  (forward-cone-only support). Produces an
  \(\mathcal{O}(\ell_P^2\nabla^2 A)\) correction to the coarse-grained
  gauge-field strength via the second moment of the multi-scale
  expansion (parallel to Appendix D's V1-second-moment derivation of
  viscosity, and Appendix D's V1-fourth-moment derivation of R1 in the
  velocity sector).
\item
  \textbf{Fourier transform.} The \(\nabla^2 A\) contribution at
  momentum \(k\) produces a \(-k^2\) shift in the gauge-field
  dispersion. For modes near the substrate cutoff \(k \sim 1/\ell_P\),
  this is an effective mass scale
\end{itemize}

\[m_\mathrm{eff}^2 \;\sim\; c_{V1}\,\ell_P^{-2},\]

with \(c_{V1}\) a dimensionless V1-profile coefficient (proportional to
the second-moment integral of the kernel, INHERITED at value layer). The
mass term is \emph{gauge-invariant kinetic-class} (a \(k\)-dependent
suppression of high-momentum modes), not Proca-class (no
symmetry-breaking mass insertion in the Lagrangian).

\begin{itemize}
\tightlist
\item
  \textbf{Non-Abelian self-interaction stabilization.} The quartic term
  \(\tfrac{1}{4}g^2(f^{abc}A^b A^c)^2\) in the YM Lagrangian provides a
  non-perturbative stabilization of the substrate-induced mass term
  against loop-correction remixing. This stabilization is structurally
  absent in the Abelian (\(U(1)\)) case (where \(f^{abc} = 0\) and the
  quartic term vanishes); it is precisely the non-Abelian
  self-interaction that distinguishes the gap-bearing case from the
  Abelian gapless case.
\item
  \textbf{Continuum-limit survival condition.} Under the limit
  \(\ell_P \to 0\) at fixed macroscopic flow scale, the bare mass scale
  \(m_\mathrm{eff,\,bare}^2 \sim \ell_P^{-2}\) diverges. The
  renormalized physical mass gap depends on how the V1 kernel profile
  rescales:
\end{itemize}

\[\boxed{\;c_{V1}(\ell_P)\,\ell_P^{-2} \;\longrightarrow\; m_\mathrm{phys}^2 \;>\; 0 \quad\text{as}\quad \ell_P \to 0.\;}\]

This is the \textbf{mass-gap survival condition}.

\textbf{Status.} The mechanism is FORCED at substrate level: V1 finite
width is unconditional (Theorem N1); the second-moment correction is
structural; non-Abelian self-interaction stabilization is FORCED by T17
clause C2. The continuum-limit survival of the gap at finite physical
value is INHERITED at value layer via the kernel-profile-rescaling
condition. ED's structural derivation supplies the mechanism but does
not pin the rescaling exponent.

This framing parallels NS-Smoothness's Intermediate Path C verdict on
Clay-NS: ED supplies a real Clay-relevant structural mechanism (R1
hyperviscous stabilization for NS smoothness; substrate-induced mass
scale for YM mass gap); the Clay-relevance verdict is conditional on a
value-INHERITED condition (R1 vs.~super-Burnett quantitative competition
for NS; kernel-profile-rescaling exponent for YM).

\begin{center}\rule{0.5\linewidth}{0.5pt}\end{center}

\subsection{E.4 Architectural Classification
(YM-4)}\label{e.4-architectural-classification-ym-4}

Following the NS-MHD-4 / Appendix C template, the six content channels
of the substrate-derived YM equation are classified under the ED-I-06
ontology:

{\def\LTcaptype{none} % do not increment counter
\begin{longtable}[]{@{}
  >{\raggedright\arraybackslash}p{(\linewidth - 6\tabcolsep) * \real{0.2500}}
  >{\raggedright\arraybackslash}p{(\linewidth - 6\tabcolsep) * \real{0.2500}}
  >{\raggedright\arraybackslash}p{(\linewidth - 6\tabcolsep) * \real{0.2500}}
  >{\raggedright\arraybackslash}p{(\linewidth - 6\tabcolsep) * \real{0.2500}}@{}}
\toprule\noalign{}
\begin{minipage}[b]{\linewidth}\raggedright
Channel
\end{minipage} & \begin{minipage}[b]{\linewidth}\raggedright
Term
\end{minipage} & \begin{minipage}[b]{\linewidth}\raggedright
ED-I-06 Class
\end{minipage} & \begin{minipage}[b]{\linewidth}\raggedright
Origin
\end{minipage} \\
\midrule\noalign{}
\endhead
\bottomrule\noalign{}
\endlastfoot
Kinetic & \(\partial_\mu F^{\mu\nu\,a}\) & Canonical ED & Participation
curvature of \(A_\mu^a\) via V1 \\
Non-Abelian self-interaction & \(gf^{abc}A_\mu^b F^{\mu\nu\,c}\) &
Canonical ED & T17 rule-type commutator (clause C2) \\
Matter source & \(J^{\nu\,a}\) & Canonical ED & T17 generalized minimal
coupling \\
Higher-derivative correction & \(c_{V1}\ell_P^2\nabla^2 A\) & Canonical
ED & V1 second-moment (YM analogue of R1) \\
Gauge condition & \(\partial_\mu A^{\mu\,a} = 0\) (Lorenz, etc.) &
Non-ED & Continuum-imposed constraint \\
Coordinate artifacts & Christoffel-class contributions & Non-ED &
Frame-kinematic bookkeeping \\
\end{longtable}
}

\textbf{Aggregate: 4 canonical ED / 1 continuum-imposed constraint / 1
frame-kinematic.}

\textbf{Three structural observations.}

\begin{itemize}
\tightlist
\item
  \textbf{YM dynamics are fully canonical ED} with no
  transport-kinematic obstruction class. The kinetic, self-interaction,
  matter-source, and higher-derivative terms --- the entire dynamical
  content of the YM equation --- all source forces from substrate
  participation structures (V1 directional-field curvature, T17
  commutator, T17 minimal coupling).
\item
  \textbf{Gauge fixing is non-ED bookkeeping.} Gauge-fixing conditions
  are continuum-imposed redundancy-removal devices rather than
  substrate-level structural features. They have no counterpart in the
  substrate-to-continuum derivation. T17 clause C8 supplies
  gauge-quotient identification at substrate level; continuum
  gauge-fixing devices are not part of the canonical-ED dynamical
  content.
\item
  \textbf{YM is structurally cleaner than NS / MHD.} NS momentum
  equation has 1 canonical / 3 non-canonical ratio; full MHD system has
  6 canonical / 5 non-canonical ratio (with 3 transport-kinematic non-ED
  items per Appendix C §C.4.3); YM equation has 4 canonical / 0
  transport-kinematic. The non-Abelian gauge sector is \emph{more}
  ED-architectural than either NS or MHD.
\end{itemize}

The structural distinction between gauge-field velocity-dependence and
fluid-velocity transport-kinematic dependence is preserved: the
non-Abelian self-interaction term \(gf^{abc}A_\mu^b F^{\mu\nu\,c}\)
might \emph{look} algebraically like a transport-kinematic bilinear, but
it arises from the substrate rule-type commutator (T17 clause C2), not
from coordinate-frame bookkeeping. Same algebraic shape as
advection-class terms; different structural origin; canonical-ED rather
than frame-kinematic.

\begin{center}\rule{0.5\linewidth}{0.5pt}\end{center}

\subsection{E.5 OS-Positivity and Continuum Stability
(YM-5)}\label{e.5-os-positivity-and-continuum-stability-ym-5}

The canonical-ED content of the YM equation (the four canonical-ED
channels of §E.4) is audited for Osterwalder-Schrader
reflection-positivity preservation under Euclidean continuation
\(t \to -i\tau\), \(A_0 \to iA_4\).

\textbf{Channel-by-channel positivity audit.}

{\def\LTcaptype{none} % do not increment counter
\begin{longtable}[]{@{}
  >{\raggedright\arraybackslash}p{(\linewidth - 6\tabcolsep) * \real{0.2500}}
  >{\raggedright\arraybackslash}p{(\linewidth - 6\tabcolsep) * \real{0.2500}}
  >{\raggedright\arraybackslash}p{(\linewidth - 6\tabcolsep) * \real{0.2500}}
  >{\raggedright\arraybackslash}p{(\linewidth - 6\tabcolsep) * \real{0.2500}}@{}}
\toprule\noalign{}
\begin{minipage}[b]{\linewidth}\raggedright
Channel
\end{minipage} & \begin{minipage}[b]{\linewidth}\raggedright
Action contribution
\end{minipage} & \begin{minipage}[b]{\linewidth}\raggedright
Positivity
\end{minipage} & \begin{minipage}[b]{\linewidth}\raggedright
Required condition
\end{minipage} \\
\midrule\noalign{}
\endhead
\bottomrule\noalign{}
\endlastfoot
Kinetic & \(\tfrac{1}{4}F_{\mu\nu}^a F_{\mu\nu}^a\) & Positive-definite
& Compact gauge group \\
Self-interaction & quartic \(g^2(f^{abc}A^bA^c)^2\) in \(S_E\) &
Positive & Compact gauge group \\
Source & \(J^\mu A_\mu\) & Non-negative bilinear & Matter-sector OS
positivity \\
Higher-derivative & \(\tfrac{1}{2}c_{V1}\ell_P^2\|\nabla A\|^2\) &
Non-negative & V1 positive Fourier transform \\
\end{longtable}
}

\textbf{Mechanisms.}

\begin{itemize}
\tightlist
\item
  \textbf{Kinetic-term positivity.} The trace-norm
  \(\tfrac{1}{2}\mathrm{Tr}(F_{\mu\nu}F_{\mu\nu}) \ge 0\) for compact
  gauge groups via Killing-form positivity (T17 clause C2). This is the
  same positivity that underlies lattice YM's plaquette action.
\item
  \textbf{Self-interaction positivity.} The Euclidean quartic term
  \(\tfrac{1}{4}g^2(f^{abc}A_\mu^b A_\nu^c)^2 \ge 0\) for compact \(G\)
  --- the structural reason constructive QFT requires compact groups for
  non-Abelian YM. The cubic term \(gf^{abc}\partial A\cdot AA\) is
  reflection-odd and contributes no negative-norm modes.
\item
  \textbf{Source-term positivity.} Bilinear
  \(\langle\Theta(JA)\,(JA)\rangle \ge 0\) conditional on matter-sector
  OS positivity (assumed via the closed Arc R / Arc Q work establishing
  fermionic and bosonic matter OS positivity).
\item
  \textbf{Higher-derivative-term positivity.} Integration by parts
  converts \(A(-\nabla^2)A\) to \(|\nabla A|^2\) in Euclidean signature;
  combined with \(c_{V1} > 0\) from V1 positive Fourier transform
  (Theorem N1 + Theorem 18), the action contribution is non-negative.
  This is the YM analogue of R1 positivity in NS-Smoothness §5 of the
  main paper.
\end{itemize}

\textbf{OS-positivity preservation locus (4 conditions).}

\begin{enumerate}
\def\labelenumi{\arabic{enumi}.}
\tightlist
\item
  \textbf{Compact gauge group} (INHERITED at value layer).
\item
  \textbf{Positive Fourier transform of V1} (FORCED at substrate level
  by T18 + N1).
\item
  \textbf{Kernel-profile rescaling}:
  \(c_{V1}(\ell_P)\,\ell_P^{-2} \to m_\mathrm{phys}^2 \ge 0\)
  (INHERITED, per §E.3).
\item
  \textbf{Matter-sector OS positivity} (structurally backed by closed
  T1--T18 work).
\end{enumerate}

If all four hold, the canonical-ED sector of the substrate-derived YM
theory is structurally compatible with OS positivity at the
channel-by-channel level audited.

\textbf{Honest framing.} The audit is structural-suggestive, not
constructively rigorous. Each canonical-ED content channel has the
correct sign and reflection structure to be compatible with OS
positivity, and the four conditions (1)--(4) are individually
satisfiable under the program's existing closed work. A constructive
proof at the Streater-Wightman / OS-axiom level would require rigorous
control of the gauge-fixing sector (Gribov-class obstructions), full
\(n\)-point Schwinger-function verification, and DCGT continuum-limit
convergence rigor --- all out of scope per the YM arc's non-goals.

\begin{center}\rule{0.5\linewidth}{0.5pt}\end{center}

\subsection{E.6 Clay-Relevance Statement
(YM-6)}\label{e.6-clay-relevance-statement-ym-6}

Following the honest framing of NS-Smoothness's Intermediate Path C
verdict:

\textbf{The ED substrate provides a structurally suggestive path toward
a constructively stable Yang-Mills continuum limit.}
Substrate-to-continuum mapping is well-defined; canonical-ED dynamics
are FORCED at substrate level; OS positivity is structurally compatible
at the canonical-ED-channel level under the four-condition preservation
locus.

\textbf{The mass-gap mechanism is FORCED at substrate level.} V1
finite-width second-moment + non-Abelian self-interaction stabilization
jointly produce a substrate-induced mass scale
\(m_\mathrm{eff}^2 \sim c_{V1}\ell_P^{-2}\). The mechanism \emph{exists}
unconditionally at substrate level; the physical mass gap value is
INHERITED via the kernel-profile-rescaling condition.

\textbf{Mass-gap survival in the continuum limit is conditional} on
\(c_{V1}(\ell_P)\ell_P^{-2} \to m_\mathrm{phys}^2 > 0\) as
\(\ell_P \to 0\). ED supplies the structural mechanism; the rescaling
exponent that determines whether the gap survives at finite positive
value is value-INHERITED.

\textbf{OS positivity is preserved channel-by-channel under the ED
canonical content} --- kinetic (positive-definite for compact \(G\)),
self-interaction (positive in Euclidean signature), matter source
(non-negative bilinear), higher-derivative (non-negative via V1
positivity). Combined Euclidean action bounded below +
reflection-positive at the structural-suggestive level.

\textbf{The gauge-fixing obstruction is reframed via substrate
gauge-quotient identification} (T17 clause C8). Standard constructive YM
has historically hit obstructions at OS-positivity preservation in the
gauge-fixing sector --- particularly via Gribov copies in
non-perturbative gauge-fixing regimes. ED's substrate-level
gauge-quotient identification reframes this from a
continuum-gauge-fixing technical problem to a
substrate-quotient-commutes-with-DCGT-coarse-graining structural
question. The reframing is structurally significant; the rigorous
resolution remains a separate technical question.

\textbf{No constructive proof of YM existence or mass gap is claimed.}
The Clay Yang-Mills Existence and Mass Gap problem requires constructive
proof of YM existence with mass gap on \(\mathbb{R}^4\) at the level of
Streater-Wightman / Osterwalder-Schrader axioms. The present analysis is
structural --- the ED account of YM existence and mass gap under
substrate-discrete-and-finite assumptions, with OS-positivity
preservation analyzed but not proven in full constructive rigor.

\textbf{The arc identifies the load-bearing structural conditions for a
positive Clay-relevance verdict:}

\begin{enumerate}
\def\labelenumi{\arabic{enumi}.}
\tightlist
\item
  \textbf{Compact gauge group} (INHERITED, value-layer).
\item
  \textbf{V1 positive Fourier transform} (FORCED, substrate-level).
\item
  \textbf{Kernel-profile rescaling condition} (INHERITED, value-layer).
\item
  \textbf{Matter-sector OS positivity} (structurally backed by closed
  work).
\end{enumerate}

If all four hold, the canonical-ED sector of the ED-derived YM theory is
structurally compatible with OS positivity and mass-gap survival. ED
supplies (2) FORCED at substrate level and structural backing for (4);
(1) and (3) are INHERITED at value layer.

\textbf{Explicitly: this is not a solution to the Clay Millennium
Problem.} It is a structural analysis identifying the conditions under
which the ED substrate → continuum mapping is compatible with the Clay
axioms. The contribution is the substrate framing of the YM equation,
the substrate origin of the mass-gap mechanism, the channel-by-channel
OS-positivity audit, and the precise specification of the
kernel-profile-rescaling condition as the load-bearing INHERITED
condition.

This Clay-relevance framing is parallel in form and honesty to the main
paper's Intermediate Path C verdict on Clay-NS: ED supplies a real
structural mechanism; identifies the load-bearing condition; does not
claim a Clay-problem solution.

\begin{center}\rule{0.5\linewidth}{0.5pt}\end{center}

\subsection{E.7 Integration with Main
Text}\label{e.7-integration-with-main-text}

This appendix interacts with the main paper and the previous appendices
in five specific ways:

\textbf{(a) Complements Appendices C and D.} Appendix C established the
ED-I-06 ontological reading of NS / MHD content under the canonical-ED /
non-ED boundary. Appendix D established the substrate-to-continuum
bridge (DCGT) for the canonical-ED dynamical content of NS / MHD.
Appendix E extends both to the non-Abelian gauge sector: the YM
architectural classification (§E.4) follows the Appendix C template; the
YM substrate-derivation (§E.2) follows the Appendix D template (DCGT
methodology generalized to non-Abelian gauge fields). The three
appendices together form a coherent architectural-and-substrate account
of the canonical-ED dynamical sector across NS, MHD, and YM.

\textbf{(b) Extends the NS Synthesis Paper into a unified NS / MHD / YM
structural foundation.} The paper, as of the addition of Appendix E,
presents: - Architectural classification (main paper §3--§6 + Appendix C
+ §E.4). - Substrate-to-continuum derivation (Appendix D + §E.2). -
ED-I-06 ontological reading (Appendix C + §E.4). - Mass-gap / smoothness
/ Clay-relevance verdicts (main paper Intermediate Path C for Clay-NS;
§E.3 + §E.6 for Clay-YM). - OS-positivity audit (§E.5).

The paper now covers the full closed-arc landscape of the program's
continuum-fluid-mechanics + non-Abelian-gauge work.

\textbf{(c) Positions the ED program for the ED-QFT unified overview
paper.} The unified-overview publication will combine the closed
structural-foundation work (T17 / T18 / Theorem N1 / Theorem GR1 / DCGT)
with the applied-arc work (NS / MHD / YM / substrate-gravity) into a
single program-overview publication. Appendix E supplies the YM building
block of that unified picture, parallel to Appendices C and D supplying
the NS / MHD building blocks. The NS Synthesis Paper, with all five
appendices in place, becomes the canonical closed-arc reference for the
continuum-fluid-mechanics + non-Abelian-gauge sector of the program.

\textbf{(d) Preserves all main-paper headline disclaimers.} The main
paper's honest disclaimers --- ED does not resolve whether 3D NS blows
up at finite time, does not predict critical Reynolds numbers, does not
supply a turbulence-cascade architectural template --- remain unchanged.
Appendix E adds parallel honest disclaimers for YM: ED does not solve
the Clay YM-existence-and-mass-gap problem, does not predict numerical
mass-gap values, does not derive the Standard Model gauge group, does
not couple to spacetime curvature. What ED \emph{does} deliver --- for
both fluid-mechanical and gauge sectors, under both structural and
substrate readings --- is a structurally coherent account of the
canonical-ED dynamical content with explicit identification of the
load-bearing conditions for Clay-relevance.

\textbf{(e) Maintains the form-FORCED / value-INHERITED methodology
across all sectors.} Every canonical-ED content channel in §E.2
(substrate YM equation), §E.3 (mass-gap mechanism), §E.4 (architectural
classification), and §E.5 (OS-positivity audit) is form-FORCED at
substrate level with values INHERITED at value layer per the program's
standard methodology. The structural verdicts are unconditional in form;
the numerical content (specific gauge group, gauge coupling, mass-gap
value, kernel-profile rescaling exponent) is INHERITED. This is the
honest methodological framing that the main paper and the previous
appendices have preserved, and Appendix E continues it.

\begin{center}\rule{0.5\linewidth}{0.5pt}\end{center}

\emph{Appendix E presents the Yang-Mills synthesis and Clay-relevance
statement. Substrate-derived continuum YM equation
\(D_\mu F^{\mu\nu} = J^\nu\) via DCGT generalized to non-Abelian gauge
fields + T17 generalized minimal coupling + non-Abelian rule-type
bracket structure. Substrate mass-gap mechanism via V1 second-moment
expansion + non-Abelian quartic stabilization, with continuum-limit
survival conditional on kernel-profile rescaling
\(c_{V1}(\ell_P)\ell_P^{-2} \to m_\mathrm{phys}^2 > 0\). Architectural
classification: 4 canonical-ED / 2 non-ED, no transport-kinematic
obstruction class --- structurally cleaner than NS or MHD. OS-positivity
audit channel-by-channel structural-suggestive positive conditional on
4-condition preservation locus (compact gauge group, V1 positive Fourier
transform, kernel-profile rescaling, matter-sector positivity).
Clay-relevance statement parallel in form and honesty to NS-Smoothness
Intermediate Path C: ED supplies the structural mechanism + identifies
load-bearing conditions; does not claim Clay-problem solution. The NS
Synthesis Paper, with Appendices C + D + E in place, presents a unified
architectural-and-substrate account of NS, MHD, and YM canonical-ED
dynamical content. The program is now positioned for the ED-QFT unified
overview paper as the next major publication move.}

\end{document}
